% ==========================================================
%  MASTER SUMMARY TEMPLATE — AI / ML / MATHEMATICS
%  Design Philosophy:
%    1. Readable running text is the default. Prose over boxes.
%    2. Boxes are used sparingly as visual landmarks for skimming.
%    3. Visual hierarchy: chapter > section > subsection > prose.
%    4. Every box is breakable — no orphaned whitespace at page ends.
%    5. A single consistent accent color ties all structure together.
% ==========================================================

\documentclass[11pt, oneside]{book}


% ══════════════════════════════════════════════════════════
%  PAGE GEOMETRY
% ══════════════════════════════════════════════════════════
\usepackage[
  a4paper,
  top    = 2.6cm,
  bottom = 2.1cm,
  left   = 2.5cm,
  right  = 2.5cm,
  headheight = 15pt
]{geometry}


% ══════════════════════════════════════════════════════════
%  TYPOGRAPHY
% ══════════════════════════════════════════════════════════
\usepackage[T1]{fontenc}
\usepackage[utf8]{inputenc}
\usepackage{lmodern}         % Latin Modern: improved Computer Modern
\usepackage{microtype}       % Microtypographic refinements (better justification,
                             %   fewer overfull hboxes — invaluable for math text)

% If inconsolata is not installed, simply remove this line.
% It provides a crisper monospace font for code listings.
\usepackage{inconsolata}

\setlength{\parindent}{0pt}
\setlength{\parskip}{0.62em}


% ══════════════════════════════════════════════════════════
%  MATHEMATICS
% ══════════════════════════════════════════════════════════
\usepackage{amsmath, amssymb, amsthm, mathtools}
\usepackage{bm}        % Bold math symbols: \bm{\theta}
\usepackage{bbm}       % Blackboard bold for indicator: \mathbbm{1}
\usepackage{mathrsfs}  % Script letters: \mathscr{F}, \mathscr{H}
\usepackage{nicefrac}  % Compact inline fractions: \nicefrac{1}{2}
\usepackage{cancel}    % Strike-through in math: \cancel{x}


% ══════════════════════════════════════════════════════════
%  SCIENTIFIC NOTATION
% ══════════════════════════════════════════════════════════
% siunitx — consistent typesetting of numbers and SI units.
%   Usage: \SI{1.5}{\micro\meter}  \SI{50}{\kilo\hertz}  \num{6.022e23}
\usepackage{siunitx}
\sisetup{per-mode = symbol}   % render "per unit" as e.g. m/s rather than m s^{-1}

% mhchem — chemical formulas and reaction equations.
%   Usage: \ce{H2O}   \ce{ATP -> ADP + P_i}   \ce{^{14}C}
\usepackage[version=4]{mhchem}


% ══════════════════════════════════════════════════════════
%  FIGURES, TABLES, LISTS
% ══════════════════════════════════════════════════════════
\usepackage{graphicx}
\usepackage{booktabs}         % \toprule, \midrule, \bottomrule
\usepackage{longtable}      % multi-page List of Abbreviations
\usepackage{array}            % Extended column types in tabular
\usepackage{float}            % [H] float placement
\usepackage{enumitem}

\setlist[itemize]{topsep=3pt, itemsep=2pt, leftmargin=2.2em}
\setlist[enumerate]{topsep=3pt, itemsep=2pt, leftmargin=2.4em}


% ══════════════════════════════════════════════════════════
%  ALGORITHMS / PSEUDOCODE
% ══════════════════════════════════════════════════════════
% Remove this block if you do not need pseudocode.
\usepackage[ruled, vlined, linesnumbered]{algorithm2e}
\SetAlFnt{\small}
\SetAlCapFnt{\small\bfseries}
\SetCommentSty{textit}


% ══════════════════════════════════════════════════════════
%  COLORS  (a refined, academic palette)
% ══════════════════════════════════════════════════════════
\usepackage{xcolor}

% Primary structural accent — deep navy used for headings, rules, ToC
\definecolor{accent}{HTML}{1B3A5C}

% Box accent colors: frame and background, designed to be legible but
% unobtrusive so running text remains the visual focus.
\definecolor{def@frame}{HTML}{1A5276}   \definecolor{def@bg}{HTML}{EBF5FB}
\definecolor{thm@frame}{HTML}{922B21}   \definecolor{thm@bg}{HTML}{FDEDEC}
\definecolor{ex@frame} {HTML}{1E8449}   \definecolor{ex@bg} {HTML}{EAFAF1}
\definecolor{int@frame}{HTML}{9A6A00}   \definecolor{int@bg}{HTML}{FEF9E7}
\definecolor{note@frame}{HTML}{5B2C6F}  \definecolor{note@bg}{HTML}{F5EEF8}

% Exam-critical marking. Slides flagged in the lecture with a GOLDEN STAR are
% "definitely exam relevant"; we echo that signal here with a matching gold.
\definecolor{stargold} {HTML}{F2B400}   % bright gold — the star glyph (≈ slide star)
\definecolor{star@frame}{HTML}{B8860B}  % readable dark goldenrod — heading text + box
\definecolor{star@bg}  {HTML}{FDF6E3}   % pale warm tint — box background


% ══════════════════════════════════════════════════════════
%  CODE LISTINGS
% ══════════════════════════════════════════════════════════
\usepackage{listings}
\lstset{
  basicstyle       = \ttfamily\small,
  breaklines       = true,
  numbers          = left,
  numberstyle      = \tiny\color{gray!70},
  frame            = tb,
  framerule        = 0.4pt,
  rulecolor        = \color{gray!35},
  tabsize          = 2,
  showstringspaces = false,
  keywordstyle     = \color{def@frame}\bfseries,
  commentstyle     = \color{gray!55!black}\itshape,
  stringstyle      = \color{ex@frame!90!black},
  backgroundcolor  = \color{gray!5},
  xleftmargin      = 2em,
  xrightmargin     = 0.4em,
  aboveskip        = 0.8em,
  belowskip        = 0.5em
}


% ══════════════════════════════════════════════════════════
%  CUSTOM BOXES  (tcolorbox)
%
%  KEY DESIGN DECISIONS:
%  1. enhanced   — enables advanced tcolorbox features (borderline etc.)
%  2. breakable  — box may split across a page break; eliminates the ugly
%                  white gap that appears when a long box would overflow.
%  3. Left-accent bar style — a thick colored stripe on the left margin.
%                  This is legible, modern, and unobtrusive relative to
%                  a full colored frame. The box background is a very
%                  light tint, keeping the reading experience calm.
%  4. sharp corners=west — left side is flat (squared), right is rounded.
%                  This mirrors the visual logic of the left accent bar.
%  5. parbox=false — allows display math to stretch properly inside boxes.
% ══════════════════════════════════════════════════════════
\usepackage[most]{tcolorbox}

\tcbset{
  enhanced,
  breakable,
  arc             = 2.5mm,
  boxrule         = 0.5pt,
  left            = 4.5mm,
  right           = 4mm,
  top             = 2.5mm,
  bottom          = 2.5mm,
  fonttitle       = \bfseries\small,
  sharp corners   = west,
  parbox          = false
}

\newtcolorbox{defbox}[1][]{
  title  = {Definition},
  colback = def@bg, colframe = def@frame,
  borderline west = {3.5pt}{0pt}{def@frame},
  #1
}

\newtcolorbox{thmbox}[1][]{
  title  = {Theorem / Key Formula},
  colback = thm@bg, colframe = thm@frame,
  borderline west = {3.5pt}{0pt}{thm@frame},
  #1
}

\newtcolorbox{exbox}[1][]{
  title  = {Example},
  colback = ex@bg, colframe = ex@frame,
  borderline west = {3.5pt}{0pt}{ex@frame},
  #1
}

\newtcolorbox{intbox}[1][]{
  title  = {Intuition \& Mental Model},
  colback = int@bg, colframe = int@frame,
  borderline west = {3.5pt}{0pt}{int@frame},
  #1
}

\newtcolorbox{notebox}[1][]{
  title  = {Note},
  colback = note@bg, colframe = note@frame,
  borderline west = {3.5pt}{0pt}{note@frame},
  #1
}

% ── Exam-critical (golden-star) markers ───────────────────
% A bright gold solid star, used to flag exam-critical headings and boxes.
\newcommand{\starglyph}{\textcolor{stargold}{\ensuremath{\bigstar}}}

% Flag a (sub)section title as exam-critical: a gold star in front and
% gold-coloured title text. \texorpdfstring keeps PDF bookmarks plain text.
% Usage:  \section{\starsec{Similarity Search}}
\newcommand{\starsec}[1]{%
  \texorpdfstring{\starglyph\hspace{0.4em}\textcolor{star@frame}{#1}}{#1}}

% Gold "exam-critical" callout box — used sparingly for the densest must-know
% facts behind a golden-star slide. (Boxes stay landmarks; do not overuse.)
\newtcolorbox{starbox}[1][]{
  title  = {\ensuremath{\bigstar}\, Exam-Critical},
  colback = star@bg, colframe = star@frame,
  borderline west = {3.5pt}{0pt}{stargold},
  #1
}


% ══════════════════════════════════════════════════════════
%  TABLE OF CONTENTS  (tocloft)
%
%  The [titles] option prevents tocloft from conflicting
%  with titlesec's redefinition of chapter/section titles.
% ══════════════════════════════════════════════════════════
\usepackage[titles]{tocloft}

% Chapter entries: accented, bold, with dots leading to page number
\renewcommand{\cftchapfont}     {\bfseries\color{accent}}
\renewcommand{\cftchappagefont} {\bfseries\color{accent}}
\renewcommand{\cftchapleader}   {\bfseries\color{accent!40}\cftdotfill{\cftdotsep}}
\setlength{\cftbeforechapskip} {7pt}

% Section entries: standard weight, small size
\renewcommand{\cftsecfont}      {\small}
\renewcommand{\cftsecpagefont}  {\small}
\setlength{\cftsecindent}      {1.8em}
\setlength{\cftbeforesecskip}  {1.5pt}

% Subsection entries: muted gray to visually recede
\renewcommand{\cftsubsecfont}      {\small\color{gray!65!black}}
\renewcommand{\cftsubsecpagefont}  {\small\color{gray!65!black}}
\setlength{\cftsubsecindent}      {3.4em}
\setlength{\cftbeforesubsecskip}  {0.5pt}


% ══════════════════════════════════════════════════════════
%  CHAPTER & SECTION HEADINGS  (titlesec)
%
%  Chapter: large faded chapter number floated right, title
%           left-aligned beneath it, then a double rule.
%           The oversized ghost number is a quiet nod to
%           modern textbook design without being gimmicky.
%
%  Section: accent-colored, with a thin rule beneath.
%           The rule draws the eye without shouting.
%
%  Subsection: slightly smaller accent color, no rule — the
%           rule at section level is already enough structure.
% ══════════════════════════════════════════════════════════
\usepackage{titlesec}

\titleformat{\chapter}[display]
  {\normalfont\bfseries\color{accent}\huge}
  {\hfill\fontsize{68}{68}\selectfont\color{accent!13}\thechapter}
  {-10pt}
  {}
  [{\vspace{4pt}%
    {\color{accent}\rule{\linewidth}{1.6pt}}\\[-2pt]%
    {\color{accent!35}\rule{\linewidth}{0.5pt}}}]
\titlespacing*{\chapter}{0pt}{-28pt}{22pt}

\titleformat{\section}
  {\normalfont\large\bfseries\color{accent}}
  {\thesection}
  {0.8em}
  {}
  [{\vspace{1pt}\color{accent!45}\titlerule[0.45pt]}]
\titlespacing*{\section}{0pt}{3.8ex plus 1ex minus .2ex}{2.2ex plus .2ex}

\titleformat{\subsection}
  {\normalfont\normalsize\bfseries\color{accent!80!black}}
  {\thesubsection}
  {0.7em}
  {}
\titlespacing*{\subsection}{0pt}{2.6ex plus 1ex minus .2ex}{1.4ex plus .2ex}

% Run-in paragraph label — useful for named sub-items in a derivation.
\titleformat{\paragraph}[runin]
  {\normalfont\small\bfseries}
  {}
  {0em}
  {}[.\enspace]
\titlespacing{\paragraph}{0pt}{1.2ex plus .5ex minus .2ex}{0.5em}


% ══════════════════════════════════════════════════════════
%  HEADER & FOOTER
% ══════════════════════════════════════════════════════════
\usepackage{fancyhdr}
\pagestyle{fancy}
\fancyhf{}
\fancyhead[L]{\small\color{gray!70}\nouppercase{\leftmark}}
\fancyhead[R]{\small\color{gray!70}\thepage}
\renewcommand{\headrulewidth}{0.4pt}

% Chapter-opening pages use LaTeX's 'plain' style — redefine it here
% to match the rest of the document (page number at foot, no header line).
\fancypagestyle{plain}{%
  \fancyhf{}
  \fancyfoot[C]{\small\color{gray!70}\thepage}
  \renewcommand{\headrulewidth}{0pt}%
}


% ══════════════════════════════════════════════════════════
%  LINKS & REFERENCES
% ══════════════════════════════════════════════════════════
\usepackage[hidelinks]{hyperref}
\usepackage[nameinlink, capitalize]{cleveref}


% ══════════════════════════════════════════════════════════
%  DOCUMENT METADATA  (fill in for each new subject)
% ══════════════════════════════════════════════════════════
\newcommand{\coursetitle}{AI in Life Sciences}
\newcommand{\documenttype}{Master Summary \\ Edition 2026}
\newcommand{\authorname}{Stefan Eder}
\newcommand{\basedon}{Based on Lecture \emph{AI in Life Sciences} SS26}


% ══════════════════════════════════════════════════════════
%  MATHEMATICAL MACROS
%  Organized by domain. Add your own at the end of each group.
% ══════════════════════════════════════════════════════════

% ── Number Sets & Spaces ──────────────────────────────────
\newcommand{\R}{\mathbb{R}}
\newcommand{\N}{\mathbb{N}}
\newcommand{\Z}{\mathbb{Z}}
\newcommand{\C}{\mathbb{C}}
\newcommand{\Q}{\mathbb{Q}}

% ── Probability & Statistics ──────────────────────────────
\newcommand{\E}{\mathbb{E}}              % Expectation
\newcommand{\Prob}{\mathbb{P}}           % Probability measure
\newcommand{\Var}{\mathrm{Var}}
\newcommand{\Cov}{\mathrm{Cov}}
\newcommand{\Corr}{\mathrm{Corr}}
\newcommand{\iid}{\overset{\mathrm{iid}}{\sim}}
% Indicator function. Usage: \ind_{\{x > 0\}}
\newcommand{\ind}{\mathbbm{1}}

% ── Information Theory ────────────────────────────────────
% KL divergence. Usage: \KL{q}{p}  →  D_KL(q ‖ p)
\newcommand{\KL}[2]{D_{\mathrm{KL}}\!\left(#1 \,\|\, #2\right)}
\newcommand{\Ent}{\mathrm{H}}            % Entropy:         \Ent(X)
\newcommand{\MI}{\mathrm{I}}             % Mutual information: \MI(X;Y)

% ── Calculus & Optimization ───────────────────────────────
\newcommand{\argmax}{\operatorname*{arg\,max}}
\newcommand{\argmin}{\operatorname*{arg\,min}}
\newcommand{\grad}{\nabla}               % Gradient
\newcommand{\hess}{\nabla^2}             % Hessian
\newcommand{\lapl}{\Delta}               % Laplacian
% Differential (upright d). Usage: \int f(x) \dd x
\newcommand{\dd}{\,\mathrm{d}}
% Partial derivatives. Usage: \pd{f}{x}  or  \pdd{f}{x}
\newcommand{\pd}[2]{\frac{\partial #1}{\partial #2}}
\newcommand{\pdd}[2]{\frac{\partial^2 #1}{\partial #2^2}}

% ── Norms, Abs, Inner Products ────────────────────────────
\newcommand{\norm}[1]{\left\lVert #1 \right\rVert}
\newcommand{\normsq}[1]{\left\lVert #1 \right\rVert^2}
\newcommand{\abs}[1]{\left\lvert #1 \right\rvert}
\newcommand{\inner}[2]{\left\langle #1,\, #2 \right\rangle}

% ── Linear Algebra ────────────────────────────────────────
\newcommand{\T}{^{\mathsf{T}}}           % Transpose:        \mathbf{A}\T
\newcommand{\inv}{^{-1}}                 % Inverse:          A\inv
\newcommand{\pinv}{^{\dagger}}           % Pseudo-inverse:   A\pinv
\newcommand{\tr}{\operatorname{tr}}
\newcommand{\rank}{\operatorname{rank}}
\newcommand{\diag}{\operatorname{diag}}
\newcommand{\spanop}{\operatorname{span}}
\newcommand{\eye}{\mathbf{I}}            % Identity matrix
\newcommand{\zeros}{\mathbf{0}}
\newcommand{\ones}{\mathbf{1}}

% ── ML / AI Core Notation ────────────────────────────────
\newcommand{\loss}{\mathcal{L}}          % Loss:          \loss(\params)
\newcommand{\reg}{\mathcal{R}}           % Regularizer:   \reg(\params)
\newcommand{\dataset}{\mathcal{D}}       % Dataset:       \dataset
\newcommand{\params}{\boldsymbol{\theta}}   % Parameters
\newcommand{\weights}{\mathbf{w}}
\newcommand{\xvec}{\mathbf{x}}
\newcommand{\yvec}{\mathbf{y}}
\newcommand{\zvec}{\mathbf{z}}
\newcommand{\hvec}{\mathbf{h}}           % Hidden state
\newcommand{\feature}{\boldsymbol{\phi}} % Feature map

% Activation functions
\DeclareMathOperator{\softmax}{softmax}
\DeclareMathOperator{\relu}{ReLU}
\DeclareMathOperator{\sign}{sign}
% sigmoid: just use \sigma (standard Greek letter)

% ── Complexity ────────────────────────────────────────────
\newcommand{\bigO}[1]{\mathcal{O}\!\left(#1\right)}
\newcommand{\bigOm}[1]{\Omega\!\left(#1\right)}
\newcommand{\bigTh}[1]{\Theta\!\left(#1\right)}

% ── Common Distributions ──────────────────────────────────
\newcommand{\Normal}{\mathcal{N}}
\newcommand{\Uniform}{\mathcal{U}}
\newcommand{\Bernoulli}{\mathrm{Ber}}
\newcommand{\Categorical}{\mathrm{Cat}}
\newcommand{\Dirichlet}{\mathrm{Dir}}
\newcommand{\Binomial}{\mathrm{Bin}}

% ── Miscellaneous ─────────────────────────────────────────
\newcommand{\given}{\,\vert\,}           % Conditional: \Prob(A \given B)
\newcommand{\st}{\;\text{s.t.}\;}        % "subject to"
\newcommand{\eqdef}{\coloneqq}           % Definitional equality  :=
\newcommand{\eqq}{\overset{?}{=}}        % "Does this equal?"


% ==========================================================
%  BEGIN DOCUMENT
% ==========================================================
\begin{document}

\frontmatter  % Roman-numeral page numbers for front matter


% ──────────────────────────────────────────────────────────
%  TITLE PAGE
% ──────────────────────────────────────────────────────────
\thispagestyle{empty}
\vspace*{1.8cm}

\begin{center}
  {\color{accent}\rule{\linewidth}{1.8pt}}
  \vspace{0.65cm}

  {\fontsize{28}{34}\selectfont\bfseries\color{accent}\coursetitle\par}
  \vspace{0.25cm}
  {\large\color{gray!80!black}\documenttype\par}

  \vspace{0.55cm}
  {\color{accent!35}\rule{\linewidth}{0.5pt}}
  \vspace{0.55cm}

  {\large\authorname\par}
  \vspace{0.2cm}
  {\small\color{gray!70}\basedon\par}
\end{center}

\vspace{1.6cm}

\begin{center}
\begin{minipage}{0.84\textwidth}
  \small\color{gray!80!black}
  This document is a compact summary of the course.
  It aims to cover all relevant concepts and results, but long derivations,
  detailed proofs, and worked examples are often condensed or omitted in
  favour of clarity and conciseness.
  Use it as a study aid alongside the lecture notes and slides, not as a
  replacement for them.

  \medskip
  If you spot an error or a gap, please open an issue or pull request at
  \textbf{github.com/stevetgit} --- corrections are always welcome.

  \medskip
  \color{gray!60!black}
  Good study materials should be shared, not gatekept.
  If this helped you, take it as a reason to help the next colleague you see struggling with something.
\end{minipage}
\end{center}

\vfill
{\centering\small\color{gray!65} \today\par}
\clearpage


% ──────────────────────────────────────────────────────────
%  TABLE OF CONTENTS
% ──────────────────────────────────────────────────────────
\setcounter{tocdepth}{2}  % Show: chapters, sections, subsections only

\begingroup
  \let\cleardoublepage\relax
  \let\clearpage\relax
  \tableofcontents
\endgroup

\mainmatter  % Switch to Arabic numerals


% ==========================================================
%  CHAPTER 1 — DRUG DISCOVERY & CHEMINFORMATICS BACKGROUND   (Lecture 01)
% ==========================================================
\chapter{Drug Discovery and the Cheminformatics Toolbox}
\label{chap:drugdiscovery}

This first chapter is mostly scene-setting. Before we can let a machine
learning model loose on molecules, we need to agree on three things: what
problem we are actually trying to solve (\emph{drug discovery}), how we feed a
molecule to a computer (\emph{representations}), and what kind of label we are
predicting (\emph{bioactivity}). None of this is deep learning yet --- but it is
the vocabulary the rest of the course is built on, so it is worth getting
comfortable with it.

The recurring idea is captured in a single picture: a function
$f_{\params}(\text{molecule}) = \text{activity}$. Everything we do later is just
a more elaborate answer to ``what is $f$, and how do we fit its parameters
$\params$ from data?''

\section{\starsec{What Drug Discovery Actually Is}}

At its core, drug discovery is the search for a molecule that \textbf{binds to a
malfunctioning biological target} and restores its function --- while staying
\emph{safe} and being \emph{processable by the body}. That last part matters as
much as the binding: a molecule that switches off the right protein but is toxic,
or that the body cannot absorb or eliminate, is not a drug.

The catch is the price tag. Bringing one drug to market takes roughly
\textbf{ten years and on the order of a billion dollars}. That is the whole
motivation for the course: if computational methods can shave even a slice off
that time and cost, the payoff is enormous.

\subsection{The Development Pipeline}

A drug does not go straight from idea to pharmacy. It moves through a pipeline,
and each stage filters out candidates:

\begin{enumerate}
  \item \textbf{Target identification} --- find the protein (or other biomolecule)
        whose malfunction drives the disease.
  \item \textbf{Drug discovery} --- find active compounds (\emph{hits} / \emph{leads}),
        optimise them into drug candidates, and worry about manufacturability
        (CMC, ``chemistry, manufacturing and controls'').
  \item \textbf{Preclinical studies} (in animals) --- safety, pharmacokinetics,
        pharmacodynamics.
  \item \textbf{Clinical trials} (in humans) --- the long, expensive phase.
  \item \textbf{Approval / launch}.
  \item \textbf{Post-market surveillance} --- watch for rare side effects once
        millions of people take it.
\end{enumerate}

AI is creeping into every stage (multi-omics analysis and knowledge graphs for
target ID, adverse-event monitoring post-market, and so on), but
\textbf{this course lives almost entirely in stage two}: finding and optimising
active small molecules.

\subsection{A Bit of Pharmacology Vocabulary}

Pharmacology is the science of how drugs interact with biological systems. Almost
every drug works by binding to a \textbf{target} --- a biological macromolecule,
usually a protein (an enzyme, receptor, ion channel, \dots), whose activity the
drug is meant to change. A molecule that binds such a target is called a
\textbf{ligand}: it slots into a specific \textbf{binding site} (or ``pocket'') on
the target and, by occupying it, modulates the target's function. ``Ligand'' is the
word we will use throughout for the small molecules we model --- and in the
ligand-based setting introduced below, the \emph{known ligands} of a target are
literally our training data. Two further words get used constantly and are easy to
mix up:

\begin{itemize}
  \item \textbf{Pharmacodynamics (PD)} --- what the \emph{drug does to the body}
        (its mechanism and effect).
  \item \textbf{Pharmacokinetics (PK)} --- what the \emph{body does to the drug}
        (how it is absorbed, distributed, metabolised, excreted).
\end{itemize}

\noindent Add \textbf{toxicology} (what harm a chemical can do) and
\textbf{drug--drug interactions} (how drugs affect each other), and you have the
main branches. Drug discovery is really the act of juggling all of these at once.

\subsection{Two Ways to Find a Drug, and Two Ways to Design One}

There is a useful historical split in \emph{how} you go hunting:

\begin{itemize}
  \item \textbf{Phenotypic} (classical / forward pharmacology): screen compounds
        for a desired \emph{effect} first, identify which lead compounds work,
        and only \emph{afterwards} figure out which target they hit.
  \item \textbf{Rational} (reverse pharmacology): start from a \emph{known target},
        screen compounds for specific interactions with it, then validate the
        effect in vivo.
\end{itemize}

Once you start using computers, a parallel split appears in how you \emph{design}
candidates --- this is \textbf{computer-aided drug design (CADD)}:

\begin{itemize}
  \item \textbf{Structure-based drug design (SBDD)} uses the \emph{3D structure}
        of the target to design molecules that fit snugly into its binding site.
        Classic tools: molecular docking and simulation. (This is the domain of
        structural biology --- not our focus here.)
  \item \textbf{Ligand-based drug design (LBDD)} uses \emph{known ligands} to
        design new molecules with similar properties, \emph{without} needing the
        protein structure. Its toolkit is chemoinformatics: QSAR, pharmacophore
        modelling, and so on.
\end{itemize}

\begin{intbox}
This whole course is essentially the \textbf{ligand-based} story. We will mostly
pretend we do \emph{not} have the protein's 3D structure and instead learn
patterns purely from molecules and their measured activities. That is exactly the
regime where machine learning on molecular representations shines --- and where
data is plentiful but messy.
\end{intbox}

\subsection{\starsec{The Design--Make--Test--Analyse Loop}}

Drug discovery is iterative. People talk about the \textbf{DMTA cycle}:
\emph{design} a molecule, \emph{make} (synthesise) it, \emph{test} it in an assay,
\emph{analyse} the result, and feed what you learned into the next design. AI now
touches each spoke of the wheel --- generative models and matched molecular pairs
for design, computer-aided synthesis planning for make, in-silico prediction for
test, and sparse-data methods for analyse. Speeding up this loop is the practical
goal behind almost everything that follows. Two things the exam likes to invert:
\emph{design comes before make and test} (you cannot test a compound you have not
yet designed and synthesised), and it is a genuine \emph{iterative cycle} --- one
round of testing never removes the need for further design.

\section{\starsec{Small Molecules and How We Represent Them}}
A quick floor of chemistry first. \textbf{Atoms} consist of protons and neutrons
in a nucleus surrounded by electrons in shells; the \emph{number of protons} fixes
the element and orders the periodic table. Only a handful of elements matter for us
--- carbon, hydrogen, oxygen and nitrogen make up the vast majority of the human
body, with phosphorus, sulfur, and a few ions following. Bonding is driven by the
outermost (\emph{valence}) electrons, the ones drawn in Lewis ``electron-dot''
symbols.

\medskip
With that in place, a drug-like small molecule is mostly a scaffold of
\textbf{carbon--carbon} and \textbf{carbon--hydrogen} bonds, decorated with a few
heteroatoms (oxygen, nitrogen, sulfur, halogens). Atoms are held together by
chemical \emph{bonds} --- chiefly \emph{covalent} (shared electrons: single,
double, triple, or delocalised ``aromatic''), but also \emph{ionic} (electrons
transferred rather than shared, as in \ce{NaCl}). Whole molecules then attract
each other through weaker \emph{intermolecular forces} (hydrogen bonding, van der
Waals, ion--dipole, \dots). These weak forces are not a footnote: they are
largely \emph{how a drug grips its target}.

A few structural subtleties matter for activity:

\begin{itemize}
  \item \textbf{3D conformation.} The same molecule can fold into different shapes
        depending on which bonds can rotate, whether rings are planar, and how it
        interacts with its environment. Shape drives binding.
  \item \textbf{Isomers.} Molecules with the same formula but different atom
        arrangements. They split into \emph{constitutional} (structural) isomers
        --- different connectivity --- and \emph{stereoisomers} --- same
        connectivity, different spatial arrangement. Stereoisomers in turn are
        either \emph{enantiomers} (non-superimposable mirror images) or
        \emph{diastereomers} (everything else), the latter including
        \emph{cis/trans} isomers and, as freely interconverting forms,
        \emph{conformers} / \emph{rotamers}.
  \item \textbf{Chirality.} A molecule is \emph{chiral} if it is not identical to
        its mirror image, like a left and right hand.
\end{itemize}

\begin{notebox}[title={Why chirality is not a technicality}]
The two mirror-image forms (\emph{enantiomers}) of a chiral drug can behave
completely differently in the body. The textbook cautionary tale is
\textbf{thalidomide}: one enantiomer is therapeutic, the other is teratogenic.
Two molecules that look almost identical on paper, wildly different effects ---
a reminder that the \emph{representation} you choose had better be able to tell
them apart.
\end{notebox}

\subsection{\starsec{``This Is Not a Molecule''}}

There is a famous slide riffing on Magritte: a drawing of caffeine captioned
\emph{``Ceci n'est pas une molecule.''} The point is genuinely important. A
SMILES string, a graph, a bit-vector --- none of these \emph{is} the molecule.
They are all lossy \emph{representations}, and the choice of representation
quietly decides what your model can and cannot learn. Picking a representation is
the first modelling decision you make, long before you choose an algorithm.

\begin{intbox}
There is no single ``true'' representation. Each one throws away some information
and keeps other information handy. A fingerprint makes substructure comparison
cheap but forgets 3D shape; a 3D geometry keeps shape but is expensive and
conformer-dependent. The art is matching the representation to the question (and
to the ML method that consumes it).
\end{intbox}

\subsection{\starsec{The Main Representations}}

\paragraph{Molecular graphs.} The most natural view: atoms are
\emph{labelled vertices}, bonds are \emph{labelled edges}, and the graph is
undirected (with some chemical constraints on edges). This is the representation
that graph neural networks consume directly (Chapter on advanced methods).

\paragraph{Molecular descriptors.} Hand-engineered numbers summarising a
molecule. They are often grouped as constitutional (molecular weight, number of
rings, count of hydrogen-bond donors/acceptors), thermodynamic (the
octanol--water partition coefficient $\log P$, torsion energy), electronic
(topological polar surface area, TPSA), steric (molecular volume), and
topological (connectivity indices). People also classify them by ``dimension''
(0D--4D) depending on how much structure they encode.

\paragraph{Molecular fingerprints.} A fingerprint is a \emph{set of features}
extracted from the molecule, then usually \emph{hashed into a fixed-length
bit-vector}. Extended-connectivity fingerprints (ECFP / Morgan) are the workhorse:
for each atom you look at growing circular neighbourhoods (diameter 0, 2, 4,
\dots), record the local environment, and hash it into the vector. Hashing is
cheap but causes \emph{bit collisions} (two different features landing on the same
bit) --- a small price for fixed-length vectors you can compare in microseconds.
ECFP is only the most common flavour; others extract different feature sets ---
\emph{atom pairs} and \emph{atom triplets} (encoding pairs/triples of atoms with
their topological distances) or \emph{(random) walks} up to a fixed length.
Conversion runs one way: a molecule (e.g.\ parsed from a \texttt{SMILES} string
into its graph) can always be reduced to a fingerprint, but the hashing is lossy
and many-to-one, so the original graph --- let alone a \texttt{SMILES} --- cannot
be reconstructed from it. Hence \texttt{SMILES} $\to$ graph $\to$ fingerprint is a
standard cheminformatics pipeline, whereas fingerprint $\to$ \texttt{SMILES} does
not exist.

\newpage

\paragraph{String representations.} Several text encodings exist; the important
ones are:

\begin{itemize}
  \item \textbf{IUPAC names} --- the human-readable systematic nomenclature
        (e.g.\ the chemical name a chemist would write). Unambiguous to a human
        but awkward for computation, which is why line notations below exist.
  \item \textbf{SMILES} (\emph{Simplified Molecular Input Line Entry
        System}) --- the dominant one. Atoms are letters (uppercase
        aliphatic \texttt{C}, lowercase aromatic \texttt{c}); bonds are implicit
        single, \texttt{=} double, \texttt{\#} triple; rings use matching digits;
        branches use parentheses; chirality uses \texttt{@}/\texttt{@@}. It is
        compact and ubiquitous --- but has \emph{strict syntax}, allows
        \emph{semantically invalid} strings, and is \textbf{non-unique} (one
        molecule has many valid SMILES).
  \item \textbf{DeepSMILES} --- a SMILES variant that reworks ring-closure and
        branch syntax to remove some classes of syntactic error, an intermediate
        step on the road from SMILES to SELFIES.
  \item \textbf{SELFIES} (\emph{Self-referencing embedded strings}) --- designed
        to fix SMILES' fragility: every SELFIES string is a \emph{100\% valid}
        molecule, which is a big deal for generative models.
  \item \textbf{SAFE} --- a SMILES-like, fragment-based encoding (fragments
        separated by dots), handy for fragment-based design.
  \item \textbf{SMARTS} --- not for whole molecules but for \emph{substructure
        patterns} (with wildcards and logical operators); it powers substructure
        search and feature extraction (e.g.\ the MACCS keys).
  \item \textbf{InChI / InChIKey} --- built for \emph{database lookup}. Because
        SMILES is non-unique, InChI normalises, canonicalises and serialises a
        molecule into a \emph{unique} string; the InChIKey is its hashed,
        fixed-length form.
\end{itemize}

\begin{notebox}[title={SMILES is non-unique --- mind the gap}]
Because a single molecule has many valid SMILES, you cannot compare molecules by
comparing SMILES strings, and naive train/test splits can accidentally place the
``same'' molecule on both sides under two different spellings. Canonicalisation
(or InChIKeys) is the fix. This non-uniqueness is exactly why InChI exists.
\end{notebox}

\begin{starbox}[title={Exam-critical: representations --- what's reversible, what isn't}]
The most-tested cheminformatics idea. \textbf{SMILES is not unique} --- one molecule
has many valid SMILES; only a \emph{canonical} SMILES (fixed algorithm) gives a
single representative string. \textbf{Direction matters:} the pipeline
SMILES~$\to$~molecular graph~$\to$~fingerprint runs \emph{one way}. You can parse a
SMILES back into a graph (that round-trip is fine), but a \textbf{fingerprint is
lossy and hashed} (bit collisions), so you \emph{cannot} reconstruct a unique graph
or SMILES from it --- ``fingerprint~$\to$~SMILES'' and ``a graph can always be
rebuilt from a fingerprint'' are both wrong. Fingerprints are \textbf{binary/count
vectors of 2D substructure}, used for similarity search and ML features --- they are
\emph{not} a 3D representation. SMILES syntax trap: \emph{lowercase} letters denote
\emph{aromatic} atoms, and malformed tokens like \texttt{[=O]} simply do not parse.
\end{starbox}
\newpage
\begin{exbox}[title={Writing a SMILES by hand: structure $\to$ string}]
A short procedure for transcribing a drawn skeleton:
\begin{enumerate}
  \item \textbf{Drop the hydrogens} --- they are implicit on C, N, O, \dots; keep
        only the heavy atoms and their bonds.
  \item \textbf{Walk a path} through the molecule from a starting atom (often a
        chain end), writing each atom's letter: uppercase aliphatic
        (\texttt{C}, \texttt{N}, \texttt{O}), \emph{lowercase} aromatic
        (\texttt{c}, \texttt{n}, \texttt{o}).
  \item \textbf{Insert bond symbols} between consecutive atoms: single is implicit,
        \texttt{=} double, \texttt{\#} triple.
  \item \textbf{Branches} --- when the path forks, write the side chain in
        parentheses \texttt{(\dots)}, then continue the main path.
  \item \textbf{Rings} --- break one ring bond and put the \emph{same digit} after
        the two atoms that were joined (a ring-closure label).
  \item \textbf{Brackets} \texttt{[\dots]} hold anything non-default: charges,
        isotopes, or explicit chirality (\texttt{[NH4+]}, \texttt{[C@H]}).
\end{enumerate}
\textbf{Worked example --- pyruvic acid} ($\text{CH}_3\text{--C}(\text{=O})\text{--C}(\text{=O})\text{--OH}$):
methyl carbon \texttt{C}; a carbonyl carbon with its double-bonded oxygen as a
branch \texttt{C(=O)}; a second carbonyl carbon \texttt{C(=O)}; the acid hydroxyl
oxygen \texttt{O} (its H implicit) $\Rightarrow$ \texttt{CC(=O)C(=O)O}.
\textbf{Rings \& aromaticity:} benzene is six aromatic carbons closed with a shared
digit, \texttt{c1ccccc1}; aspirin combines all the rules ---
\texttt{O=C(C)Oc1ccccc1C(=O)O} (acetyl ester on an aromatic ring bearing a
carboxyl). Any valid spelling parses to the same molecule; a \emph{canonical}
SMILES just fixes one.
\end{exbox}

\paragraph{Geometry- and physics-based representations.} A further family keeps
the molecule's \emph{3D structure} or its underlying \emph{physics}. These show up
mainly in quantum- and property-prediction (QSPR) rather than ligand-based QSAR:

\begin{itemize}
  \item \textbf{3D geometry} --- each atom as a tuple $(Z_i, x_i, y_i, z_i)$:
        nuclear charge plus Cartesian coordinates. Physically complete, but
        \emph{conformer-dependent} and not invariant to rotation/translation.
  \item \textbf{Coulomb matrix} --- a fixed-size square matrix encoding nuclear
        Coulomb repulsion, off-diagonal
        $M_{ij}\propto Z_iZ_j/\lVert\mathbf{r}_i-\mathbf{r}_j\rVert$ and diagonal a
        fit to free-atom energies (defined only up to atom ordering).
  \item \textbf{Bag of bonds / fragments} --- counts (a ``bag'') of bond types or
        substructure fragments, discarding their global arrangement.
  \item \textbf{Potentials} --- molecular interaction / electrostatic potentials
        sampled in the space around the molecule.
  \item \textbf{Electronic density} $\Psi$ --- the wavefunction / electron density:
        the most information-rich but by far the most expensive to obtain.
\end{itemize}

\begin{intbox}
The 2D representations above (graphs, fingerprints, strings) dominate
\emph{ligand-based} work: they are cheap and need no 3D conformer. The
geometry/physics representations carry more information but cost more and hinge on
a chosen conformation --- which is why they live mostly in quantum-property
prediction, not in the data-hungry QSAR this course is about.
\end{intbox}

Finally, modern pipelines increasingly use \textbf{learned} (continuous)
representations produced by deep networks, and \textbf{biological} descriptors
such as compound-induced transcriptomics, interaction fingerprints, or
pharmacophores --- representations that try to capture \emph{biological} rather
than purely \emph{chemical} similarity. We will come back to that distinction
in Chapter~\ref{chap:vsqsar}.

\section{\starsec{Biological Activity and Assays}}
We already carry a lot of informal biological knowledge about small molecules ---
sugars taste sweet, ibuprofen relieves pain, caffeine is a stimulant. Historically
that knowledge accrued in experts' heads; today we acquire and store it
\emph{systematically}, as measured numbers attached to (compound, target) pairs.
The unit of that measurement is the assay.

\textbf{Bioactivity} is a compound's ability to interact with a biological system
and produce a \emph{measurable response}. We find out about it by running a
\textbf{bioassay} --- a systematic testing procedure that measures the effect of a
compound on a biological system and returns a numerical or qualitative outcome.
Every bioassay has three parts worth naming:

\begin{itemize}
  \item the \textbf{substance} being tested (e.g.\ a drug candidate),
  \item the \textbf{target} / biological system (a protein, cell, or organism),
  \item the \textbf{endpoint}, i.e.\ the measured number (IC$_{50}$, $K_d$,
        \% inhibition, \dots).
\end{itemize}

Like any measurement, assays produce \emph{false positives, false negatives and
biases} --- a theme that will haunt us when we evaluate models. Assays come in
flavours: \emph{in vitro} (cell-free or cell-based), \emph{ex vivo}
(tissue-based), and \emph{in vivo} (whole organism); they can measure
\emph{binding} (affinity to a target) or \emph{function} (an actual biological
response); and \emph{high-throughput screening (HTS)} automates testing on a large
scale.

\begin{exbox}[title={An \emph{in vitro} assay: the Ames test}]
A classic cell-based example is the \textbf{Ames test} for mutagenicity. A
histidine-dependent \emph{Salmonella} strain is plated with the test compound (and
rat-liver extract to mimic metabolism) on histidine-poor medium. Only
back-mutated (\emph{his$^-$}$\to$\emph{his$^+$}) revertants can grow; a colony
count well above the natural-revertant control flags the compound as a likely
mutagen. The endpoint here is \emph{qualitative-to-counted}, not an XC$_{50}$.
\end{exbox}

\begin{starbox}[title={Exam-critical: \emph{in vitro} vs.\ \emph{in vivo}}]
A definitional trap. \textbf{\emph{In vitro}} means ``in glass'' --- a test tube,
plate, or cell culture (cell-free \emph{or} cell-based); it does \emph{not} mean
``in a living organism''. \textbf{\emph{In vivo}} is the whole living organism, and
\textbf{\emph{ex vivo}} is excised tissue. So an option that calls \emph{in vitro}
``testing in a living organism'' is wrong (that is \emph{in vivo}). Otherwise,
assays do measure a biological response/function, are run over a concentration
range (dose--response), and are physical wet-lab experiments --- those statements
are all correct.
\end{starbox}
\newpage
\noindent Which branch of pharmacology you are probing dictates the assay you run
and the number it returns:

\begin{table}[H]
\centering
\small
\begin{tabular}{lll}
\toprule
\textbf{Branch} & \textbf{Example assay} & \textbf{Typical output(s)} \\
\midrule
Pharmacodynamics        & Receptor binding assay         & Binding affinity $K_i$, IC$_{50}$ \\
Pharmacokinetics        & Microsomal stability assay     & Metabolic half-life, clearance rate \\
Toxicology              & Cytotoxicity (cell viability)  & CC$_{50}$ (50\% cell death), LD$_{50}$ (lethal dose) \\
Drug--drug interactions & CYP450 inhibition assay        & IC$_{50}$ for enzyme inhibition \\
\bottomrule
\end{tabular}
\caption{Representative assays and their endpoints across the main pharmacology
branches. The XC$_{50}$ family (next subsection) generalises IC$_{50}$,
CC$_{50}$, and related half-maximal measures.}
\end{table}

\subsection{Dose--Response Curves and the \texorpdfstring{XC$_{50}$}{XC50} Family}

A compound's effect depends on its \emph{concentration} (in mol/L, usually
reported in nanomolar). If you test a substance at many concentrations you trace
out a \textbf{dose--response curve}, assumed to be sigmoidal. The single number
people quote is the concentration producing \emph{half} the maximal effect,
generically \textbf{XC$_{50}$} --- e.g.\ \textbf{IC$_{50}$} is the concentration
causing 50\% inhibition.

Because these concentrations span many orders of magnitude, we put them on a log
scale. The convention (with concentration $X$ in molar) is
\[
  \mathrm{pX} \;=\; -\log_{10} X .
\]

\begin{thmbox}[title={Activity on a log scale}]
With $X$ in molar, $\mathrm{pX} = -\log_{10} X$. Hence
\[
  \text{lower XC}_{50} \;\Longleftrightarrow\; \text{higher pX}
  \;\Longleftrightarrow\; \text{higher activity.}
\]
A potent compound has a \emph{small} IC$_{50}$ but a \emph{large} pIC$_{50}$. Get
the direction wrong and your model optimises for exactly the opposite of what you
want.
\end{thmbox}

\begin{exbox}[title={From IC$_{50}$ to pIC$_{50}$}]
Suppose an assay reports $\mathrm{IC}_{50} = 28\,\text{nM} = 28\times10^{-9}\,$M.
Then
\[
  \mathrm{pIC}_{50} = -\log_{10}\!\left(2.8\times10^{-8}\right) \approx 7.55 .
\]
A second compound at $\mathrm{IC}_{50} = 2.8\,\mu$M gives
$\mathrm{pIC}_{50}\approx 5.55$ --- two log units less potent. Modelling pIC$_{50}$
rather than IC$_{50}$ keeps the target variable on a sane, roughly Gaussian scale.
\end{exbox}

Other affinity measures ($K_i$, $K_d$) play the same role. The takeaway: these
numbers are the \emph{labels} our QSAR models will try to predict.

\newpage

\subsection{\starsec{Pharmacokinetics: ADME and Drug-Likeness}}

Good binding is necessary but not sufficient. The compound also has to reach its
target and leave the body cleanly. That is summarised by \textbf{ADME}:

\begin{itemize}
  \item \textbf{A}bsorption --- how it enters the bloodstream,
  \item \textbf{D}istribution --- how it spreads to tissues (e.g.\ crossing the
        blood--brain barrier),
  \item \textbf{M}etabolism --- how it is chemically modified,
  \item \textbf{E}xcretion --- how it is cleared.
\end{itemize}

ADME is genuinely hard to model, so people lean on cheap proxies for
``drug-likeness''. The classic is \textbf{Lipinski's Rule of Five}.

\begin{defbox}[title={Lipinski's Rule of Five (Ro5)}]
A compound is more likely to be orally bioavailable if it has
\emph{no more than} 5 hydrogen-bond donors, \emph{no more than} 10 hydrogen-bond
acceptors, molecular weight $< 500$, and $\log P < 5$. The
\textbf{Quantitative Estimate of Drug-likeness (QED)} is a softer, continuous
alternative: it scores how closely a molecule's \emph{eight} physicochemical descriptors resemble those of
$\sim$771 approved oral drugs.
\end{defbox}

\begin{starbox}[title={Exam-critical: what good ADME does (and doesn't) buy}]
\textbf{ADME = pharmacokinetics} (what the body does to the drug over time), and
Lipinski's Ro5 is a cheap drug-likeness \emph{proxy} for it --- both true. The trap
is sufficiency: good ADME is \textbf{necessary but not sufficient} for an effective
drug. It only means the compound can reach its target intact; efficacy still needs
actual target \emph{binding} and the right pharmacodynamics. Also true and easily
missed: ADME endpoints such as \emph{clearance} depend on complex metabolism and are
notoriously hard to model.
\end{starbox}

\subsection{\starsec{Why the Data Is Painful}}

A theme that will not go away: bioassay data is \textbf{limited, imbalanced and
noisy}.

\begin{itemize}
  \item \textbf{Expensive labels.} A single high-quality data point can cost up to
        \$10{,}000 to measure.
  \item \textbf{Small datasets.} Many assays cover just 20--500 molecules; HTS
        campaigns reach 10k--100k but are noisier.
  \item \textbf{Imbalance.} Almost everything is \emph{inactive}; actives are rare.
  \item \textbf{Noise.} Typical measurement noise is $\sim 0.5$ log units, and it
        \emph{grows} when you pool data across different assays.
\end{itemize}

\begin{notebox}[title={Set expectations accordingly}]
Limited + imbalanced + noisy is the default, not the exception. It is why so much
of QSAR is about \emph{data-efficient, robust} learning and careful evaluation,
rather than chasing the fanciest architecture. Keep this list in mind --- it
justifies almost every design choice in Chapter~\ref{chap:vsqsar}.
\end{notebox}

\section{\starsec{Where the Data Lives: Databases and Chemical Space}}

The space of possible drug-like molecules is staggeringly large, and the part we
have actually measured is a tiny sliver of it.

\begin{table}[H]
\centering
\small
\begin{tabular}{lr}
\toprule
\textbf{Region of chemical space} & \textbf{Approx.\ size} \\
\midrule
Estimated drug-like compounds (Lipinski) & $10^{33}$--$10^{60}$ \\
Enumerated virtual library (GDB-17)       & $\sim 10^{11}$ \\
On-demand commercial library (ZINC-22)    & $\sim 10^{10}$ \\
Synthesised / in-stock (PubChem, ZINC20)  & $\sim 10^{8}$ \\
Approved drugs (DrugBank)                  & $\sim 10^{3}$ \\
\bottomrule
\end{tabular}
\caption{Orders of magnitude in chemical space. We have made $\sim 10^{8}$
molecules and approved $\sim 10^{3}$ --- out of perhaps $10^{33}$ or more.}
\end{table}

For \emph{bioactivity} data specifically there are two camps:

\begin{itemize}
  \item \textbf{Structure-based data} (PDB, PDBBind; $> 10$k points): you get the
        binding site and 3D structure, but there are \emph{no negatives}
        (inactives) and only a small number of ligands.
  \item \textbf{Ligand-based data} (PubChem, ChEMBL; $> 100$M points): huge, with
        \emph{labelled actives and inactives} --- but no binding conformation.
\end{itemize}

Since we are doing ligand-based modelling, the second camp is home. The main
public databases:

\begin{table}[H]
\centering
\small
\begin{tabular}{lrr}
\toprule
\textbf{Database} & \textbf{\#Compounds} & \textbf{\#Bioassay datapoints} \\
\midrule
PubChem   & 117\,M & 294\,M \\
ChEMBL    & 2.5\,M & 21\,M \\
Papyrus   & 1.2\,M & 60\,M \\
BindingDB & 1.3\,M & 2.9\,M \\
DrugBank  & 20\,K  & --- \\
\bottomrule
\end{tabular}
\caption{Public ligand-based bioactivity databases. A ``datapoint'' is a triple
(molecule, assay, value). A lot of additional data sits in private company
archives or in non-electronic literature.}
\end{table}

\begin{notebox}[title={The data is mostly holes}]
Take ChEMBL: $\sim$2.5M compounds, $\sim$16k targets and $\sim$1.7M distinct
assays, but only $\sim$21M activity values --- on the order of $\sim$12 datapoints
per assay and $\sim$8 per compound. If you arrange compounds against assays as a
matrix, roughly \textbf{0.0005\%} of the cells are filled. QSAR modelling is, to a
large extent, the problem of doing something useful with an almost-empty matrix. 
\end{notebox}
\newpage
\section{\starsec{From SAR to QSAR (and QSPR)}}

The empirical bedrock under all of this is the \textbf{structure--activity
relationship (SAR)}: a molecule's biological activity is determined by its
structure, and \emph{similar structures tend to have similar activities}. The
\emph{quantitative} version, \textbf{QSAR}, claims there is an underlying
\emph{mathematical} relationship between structure and bioactivity --- precisely
our $f_{\params}(\text{molecule}) = \text{activity}$.

\begin{defbox}[title={QSAR / QSPR}]
\textbf{QSAR} (\emph{quantitative structure--activity relationship}) models the
\emph{biological} activity of a compound from its structure.
\textbf{QSPR} (\emph{quantitative structure--property relationship}) does the same
for physico-chemical \emph{properties} (solubility, melting point, chemical
stability). QSPR data is usually easier and cheaper to obtain, because it does
not involve a living system.
\end{defbox}

\begin{starbox}[title={Exam-critical: what QSAR is --- and isn't}]
\textbf{Scope:} QSAR maps \emph{molecule~$\to$~property} (activity, toxicity).
Predicting 3D \emph{protein} structure (e.g.\ AlphaFold) is \emph{structure
prediction}, and genome mapping is genomics --- neither is QSAR, even though both
are ``AI in life science''. \textbf{Data shape:} typical bioactivity datasets are
\textbf{limited, imbalanced and noisy} --- hundreds to a few thousand points per
target (\emph{not} millions), far more inactives than actives (\emph{not} balanced),
with measurement noise, and heavily clustered into \emph{analog series} (one
scaffold, small variations). Any option calling the data ``large, balanced and
clean'' is the reverse of reality.
\end{starbox}

QSAR is older than machine learning by a century. A quick lineage: Crum-Brown \&
Fraser (1868) proposed that physiological action depends on chemical composition;
Richet (1893) linked cytotoxicity inversely to water solubility; Meyer and Overton
(\textasciitilde 1900) tied narcotic action to oil/water partitioning; and
Hansch \& Muir (1962) related activity to Hammett constants and hydrophobicity ---
widely called the \emph{birth of QSAR}.


The \emph{modelling} methods then evolved roughly as:

\begin{center}
\small
\begin{tabular}{ll}
\toprule
\textbf{Era} & \textbf{Dominant methods} \\
\midrule
$<$ 1990   & linear models, simple heuristic relationships \\
$\sim$1990s & artificial neural networks \\
$\sim$2000s & SVMs, random forests, gradient boosting \\
2014--      & deep nets, multi-task nets, graph neural nets \\
\bottomrule
\end{tabular}
\end{center}

\begin{intbox}
Notice the through-line: \emph{represent} the molecule as numbers, then
\emph{regress or classify} activity from those numbers. The decades just changed
\emph{how} we get the numbers and \emph{how flexible} the function $f$ is. With
that framing in hand, the next chapter is simply: how do we \emph{use} such a
model to screen a library, and how do we \emph{build and trust} one in the first
place (both halves of Chapter~\ref{chap:vsqsar})?
\end{intbox}


% ==========================================================
%  CHAPTER 2 — VIRTUAL SCREENING & QSAR MODELING   (Lecture 02)
% ==========================================================
\chapter{Virtual Screening \& QSAR Modeling}
\label{chap:vsqsar}

We now have molecules, representations, and a notion of activity. The first thing
we want to \emph{do} with all that is sift an enormous pile of candidate molecules
down to the handful worth testing in the lab. That is \textbf{virtual screening
(VS)}, and it is the central task of computer-aided drug design. The natural
second question --- how do we \emph{build} the model that does the sifting, and
how do we know whether to trust it? --- is \textbf{QSAR modelling}, the second
half of this chapter. The two belong together: VS is what we want to do, QSAR is
how we learn to do it well.

\section{\starsec{The Idea: Screening a Library In Silico}}

Picture a compound library of more than a million molecules. Testing all of them
in assays is expensive or outright infeasible. So we insert a computational filter:

\begin{center}
\emph{molecule library} ($>$1M) $\;\to\;$ \textbf{virtual screening}
$\;\to\;$ selected molecules $\;\to\;$ wet-lab assay.
\end{center}

The model evaluates each molecule for the desired property and \emph{prioritises}
a small, promising subset for actual experiments. The payoff is concrete: it
shrinks the search space, boosts the efficiency of HTS, and cuts cost
(HTS alone accounts for $\sim$15\% of pharmaceutical research spending).

\begin{intbox}
The mental model is a \emph{funnel}, and the model's job is \textbf{ranking}, not
perfect prediction. We do not need the exact activity of all million molecules ---
we need the truly active ones to float to the top of the list so that the limited
assay budget is spent where it counts. This ``ranking matters'' view will shape
how we evaluate models later in this chapter.
\end{intbox}

\subsection{Ligand-Based vs.\ Structure-Based VS}

There are two families of methods, mirroring the SBDD/LBDD split from
Chapter~1:

\begin{itemize}
  \item \textbf{Ligand-based} VS uses \emph{only ligand information} and leans on
        SAR. If you know only a \emph{few} actives, you do a \textbf{similarity
        search}; if you have \emph{enough} actives and inactives, you build a
        predictive \textbf{model} (model-based search).
  \item \textbf{Structure-based} VS uses \emph{both} ligand and protein structure.
        Historically this meant physics-based methods (docking, molecular dynamics,
        free-energy perturbation); increasingly it means deep-learning docking and
        co-folding.
  \item \textbf{Hybrid} approaches use both, e.g.\ proteochemometrics (PCM) models
        and contrastive methods that bridge several modalities.
\end{itemize}

\noindent The rest of this chapter is \textbf{ligand-based}, in two modes:
similarity search and model-based search.

\begin{starbox}[title={Exam-critical: four definitions students swap}]
\textbf{SBDD vs.\ LBDD:} structure-based uses the target's \emph{3D structure};
ligand-based uses \emph{known ligands without} the target structure (the trap
reverses these). Molecular fingerprints are an LBDD tool, and LBDD \emph{is} used to
predict activity of untested compounds by analogy --- selection also weighs ADME and
structural alerts, not bioactivity alone. \textbf{Similarity search vs.\ model-based
search:} similarity search \emph{retrieves neighbours} of a query molecule (good
when only a few actives are known); model-based search trains a \emph{supervised
model} on labelled actives and inactives (and, being ligand-based, needs \emph{no}
target 3D structure). The trap reverses these two definitions too.
\end{starbox}

\section{\starsec{Similarity Search}}

The premise is the SAR slogan taken literally: if a query molecule is
\emph{similar} to a known active, it is probably active too. So we rank a library
by similarity to a known active and test the top hits.

The hard part is hiding in the word ``similar''. Similarity can mean
\emph{chemical}, \emph{pharmacological}, or \emph{biological} resemblance, and
these do not coincide. Two molecules can look chemically alike yet behave
differently, or look chemically different yet hit the same target (as families of
MAPK inhibitors do).

\begin{intbox}
There is no universally ``correct'' similarity measure --- if there were, drug
design would essentially be solved. Every choice of fingerprint and metric encodes
a particular guess about \emph{what makes molecules behave alike}. Choosing one is
a modelling assumption, not a neutral preprocessing step.
\end{intbox}

\subsection{\starsec{Fingerprint Similarity}}

The cheapest and most common approach compares binary fingerprints (Morgan/ECFP,
RDKit, MACCS). Let $a$ be the number of features (set bits) in molecule $A$, $b$
the number in $B$, and $c$ the number they share. Three metrics dominate:

\begin{thmbox}[title={Fingerprint similarity coefficients}]
\begin{align*}
\text{Tanimoto (Jaccard / IoU):}\quad
  & T(A,B) = \frac{|A\cap B|}{|A\cup B|} = \frac{c}{a+b-c},\\[2pt]
\text{Dice:}\quad
  & D(A,B) = \frac{2|A\cap B|}{|A|+|B|} = \frac{2c}{a+b},\\[2pt]
\text{Cosine:}\quad
  & C(A,B) = \frac{A\cdot B}{\norm{A}\,\norm{B}} = \frac{c}{\sqrt{a\,b}}.
\end{align*}
More shared features $\Rightarrow$ higher similarity. \textbf{Tanimoto} is the
default in cheminformatics.
\end{thmbox}

\begin{exbox}[title={Computing all three}]
Take $a = 3$, $b = 5$, $c = 2$ (molecules with 3 and 5 set bits, sharing 2):
\[
  T = \frac{2}{3+5-2} = \tfrac{1}{3} \approx 0.33,\qquad
  D = \frac{2\cdot 2}{3+5} = 0.5,\qquad
  C = \frac{2}{\sqrt{15}} \approx 0.52 .
\]
Same molecules, three different numbers --- which is why ``how similar are they?''
only has an answer once you fix \emph{both} the fingerprint and the metric.
\end{exbox}

\subsection{\starsec{Beyond Fingerprints}}

Other chemical-similarity notions include the \emph{maximum common substructure},
\emph{graph edit distance}, \emph{Bemis--Murcko scaffolds} (the ring-system
backbone of a molecule), and \emph{ECFP kernels}.

A more ambitious idea is \textbf{pharmacophore similarity}. A pharmacophore is an
\emph{abstract} description of a molecule --- typically 3--4 points encoding
features like hydrogen-bond donors/acceptors, polarity, hydrophobicity and charge.
It is meant to capture \emph{biological} rather than chemical similarity, and is
usually handcrafted by experts (there is no agreed-upon definition). One can then
define a kernel $k(\cdot,\cdot)$ between the pharmacophore ``graphs'' of two
molecules.

Finally, similarity can be \textbf{learned}. Train a neural network whose input is
a \emph{chemical} representation and whose output predicts \emph{biological}
activity; the hidden layers then form a representation that has been pushed from
chemistry towards biology. Comparing molecules with a kernel
$k\big(h(A), h(B)\big)$ on a hidden layer $h$ gives a similarity that reflects
biological behaviour better than raw chemical overlap.

\begin{intbox}
The arc of this section is a slow migration from \emph{chemical} to
\emph{biological} similarity: hard-coded fingerprints $\to$ expert pharmacophores
$\to$ learned embeddings. Each step tries to make ``similar'' mean ``behaves the
same in the body'' rather than ``looks the same on paper''.
\end{intbox}

\subsection{Running the Search}

With a similarity (or kernel) $k$ in hand, the search itself is simple: given a
known active, compute $k(\text{active}, m)$ for every molecule $m$ in the library
and return the most similar ones. With a million molecules this is just a sparse
vector--matrix operation --- fast enough to do at scale.

\section{\starsec{Model-Based Search}}

When you have \emph{enough} labelled actives \emph{and} inactives, you can do
better than nearest-neighbour similarity: train a \textbf{supervised model} that
learns its own rules. Concretely, fit a function
\[
  y = g(\xvec; \weights)
\]
that maps a molecule's representation $\xvec$ either to a class (active vs
inactive) or to a continuous activity value, then apply it across the library and
rank by the predicted score. This is the bridge into QSAR modelling proper, which
fills the rest of this chapter.

\newpage

\section{\starsec{Filtering and Clustering the Hits}}

A ranked list of predicted actives is not the end. Two practical clean-up steps
make the shortlist actually useful.

\subsection{Filtering Out Problematic Structures}

Virtual screening focuses on drug--target interaction, but ADME and toxicity
matter just as much for a \emph{viable} candidate. So we drop structures that are
likely to cause trouble downstream:

\begin{itemize}
  \item \textbf{Lipinski Rule of Five} --- a coarse drug-likeness gate.
  \item \textbf{Brenk filters} --- known toxic / reactive substructures.
  \item \textbf{PAINS} (pan-assay interference compounds) --- scaffolds notorious
        for producing \emph{false positives} across many assays.
  \item \textbf{hERG inhibition} --- a flag for potential cardiotoxicity via the
        hERG ion channel.
\end{itemize}

\begin{notebox}[title={PAINS: the false-positive trap}]
PAINS are nasty precisely because they ``hit'' in lots of assays without being
real binders --- they interfere with the readout. If you do not filter them, your
model happily learns to predict assay artefacts as activity. Filtering is part of
honest virtual screening, not an afterthought.
\end{notebox}

\subsection{Clustering for Diversity}

If your top-ranked hits are all minor variations of one another, testing all of
them wastes the assay budget on redundant information. \textbf{Clustering} keeps
the shortlist \emph{diverse}, maximising the chance of finding genuinely novel
chemotypes. The recipe: compute pairwise similarity (from fingerprints or
descriptors), cluster, and pick a representative or two from each cluster to span
the chemical space.

Two broad clustering styles appear:

\begin{itemize}
  \item \textbf{Distance-based}: $k$-means, DBSCAN, hierarchical, and
        \emph{sphere-exclusion} methods (Butina, MaxMin) using a molecular distance
        (e.g.\ $1-$Tanimoto).
  \item \textbf{Graph-based}: grouping by Bemis--Murcko scaffold or maximum common
        substructure.
\end{itemize}

This idea of similarity-aware grouping comes back with a vengeance later in this
chapter, where the \emph{same} clustering is used not to diversify
hits but to build honest train/test splits.

\subsection{A Note on Tooling}

In practice this is all Python. \texttt{rdkit} is the backbone --- it converts
between molecule formats, handles databases, does substructure search via SMARTS,
and generates descriptors and fingerprints. Similarity search uses
\texttt{scipy}/\texttt{numpy}/\texttt{rdkit} for sparse vector operations; classic
ML models come from \texttt{sklearn} (random forests, SVMs); and neural networks
use \texttt{PyTorch}/\texttt{TensorFlow}/\texttt{Keras}.


% ----------------------------------------------------------
%  (former Chapter 3 — QSAR MODELING; now the second half of this chapter)
% ----------------------------------------------------------

\newpage
\section{\starsec{The QSAR Modelling Pipeline}}

So far we used a model as a black box to rank a library; we now open the box. How
do we actually \emph{build} a model
$f_{\params}(\text{molecule}) = \text{activity}$, and --- just as important --- how
do we \emph{know whether to trust it}? With molecular data, that second question
turns out to be where most of the real difficulty lives.

A QSAR project is a pipeline, and it pays to see the whole assembly line before
zooming in:

\begin{center}
bioactivity data $\;\to\;$ molecular representation $\;\to\;$ data splits
$\;\to\;$ feature selection $\;\to\;$ model \& hyperparameter selection
$\;\to\;$ training $\;\to\;$ evaluation.
\end{center}

The input is the familiar triple \emph{(assay, molecule, value)} --- and, as we
keep stressing, that data is limited, noisy, imbalanced and biased. Every box in
the pipeline is really a defence against one of those four problems.

\subsection{\starsec{Representation and Feature Selection}}

In principle \emph{any} numerical representation from Chapter~1 can feed a QSAR
model; which one works best depends on the ML method paired with it. Feature
selection can be applied as usual --- with one crucial caveat.

\begin{notebox}[title={Feature-selection bias}]
Feature selection must happen \emph{inside} the training loop, using only training
data. If you pick features by looking at the whole dataset (test set included),
information leaks from test to train and your performance estimate is optimistic.
The same discipline applies to scaling, imputation, and any other ``learned''
preprocessing.
\end{notebox}

A useful way to organise the representation zoo is by \emph{dimensionality}, since
it tends to predict which model family fits:
\begin{itemize}
  \item \textbf{0D --- descriptors}: handcrafted molecular features \emph{or}
        \textbf{NN-generated (learned)} descriptors, fed to gradient-boosted trees
        (\textbf{CatBoost}) or feed-forward DNNs.
  \item \textbf{1D --- text}: SMILES strings, fed to a \textbf{Transformer} or CNN.
  \item \textbf{2D --- graph}: the molecular graph, fed to graph networks
        (\textbf{KGCNN}-style models).
\end{itemize}
The richer representations from Chapter~1 (potentials, weighted graphs, Coulomb
matrices, bag of bonds, 3D geometry, electronic density) carry over unchanged; the
point here is only that representation and method are chosen \emph{together}.

\subsection{\starsec{Coping With Limited, Noisy, Imbalanced Data}}

Recall the pain points from Chapter~1: \$10k labels, datasets of 20--500
molecules (high-throughput campaigns reach 10k--100k, but are noisier), far more
inactives than actives, and $\sim$0.5 log-unit noise. Two
pragmatic responses recur throughout QSAR:

\begin{itemize}
  \item Use \textbf{robust, data-efficient} learning methods (the small-data
        regime rewards strong inductive biases over raw capacity).
  \item When labels are too noisy for regression, \textbf{switch to
        classification} by thresholding the activity value (active vs inactive).
        Predicting a coarse label is often more reliable than predicting a noisy
        number.
\end{itemize}

\section{\starsec{Which Model? Representations Meet Algorithms}}

The algorithmic menu tracks the historical timeline from Chapter~1: linear models,
then SVMs / random forests / gradient boosting, then deep and graph networks. In
QSAR the four everyday workhorses are \textbf{random forests}, \textbf{SVMs},
\textbf{gradient boosting}, and \textbf{neural networks}.

What matters in practice is that the \emph{representation} and the \emph{method}
have to fit each other:

\begin{table}[H]
\centering
\small
\begin{tabular}{l p{8.4cm}}
\toprule
\textbf{Representation} & \textbf{Methods it pairs well with} \\
\midrule
Descriptors / fingerprints & linear models, RF, gradient boosting, SVM (standard kernel), feed-forward NN \\
SMILES strings             & RNN/LSTM, CNN or transformer, SVM (string kernel) \\
Molecular graph            & graph neural networks, SVM (graph / molecule kernel) \\
Bag of fragments           & SVM (Tanimoto kernel), DNN with set pooling \\
\bottomrule
\end{tabular}
\caption{Common (not exhaustive) representation--method combinations. Empty cells
in the lecture's version mean ``non-trivial'', not ``impossible''.}
\end{table}

\begin{intbox}
The honest summary is: \textbf{there is no single optimal architecture}. A
gradient-boosted model on ECFP fingerprints is a famously hard baseline to beat on
small datasets, while graph networks shine when data is plentiful. Picking the
representation/method pair is itself a hyperparameter --- which is exactly why the
next two sections are about \emph{choosing} and \emph{validating} honestly.
\end{intbox}

\subsection{Model and Hyperparameter Selection}

A lucky break of QSAR is that models are often \emph{cheap} to train (small
datasets), so you can afford to try many architectures and hyperparameter
settings. The danger is the flip side: try enough configurations and one will look
good \emph{by chance}.

\begin{notebox}[title={Hyperparameter selection bias $\to$ use nested CV}]
Never tune hyperparameters on the same data you report performance on. The clean
solution is \textbf{nested cross-validation}: an \emph{inner} loop selects
hyperparameters, and a separate \emph{outer} loop estimates generalisation. Any
automated \emph{feature selection} belongs to the outer loop too --- performed
once per outer fold, on that fold's training portion only. The outer loop never
sees the choices made on its test folds, so the final number is unbiased.
\end{notebox}

\section{\starsec{The Big QSAR Pitfall: Data Splits}}

If you remember one thing from this chapter, make it this section. The way you
\emph{split} molecular data into train and test sets can quietly inflate your
reported performance by a lot.

\subsection{\starsec{Compound Series Bias}}

QSAR datasets are typically built from \emph{series} of closely related molecules
--- chemists make many small variations on a promising scaffold, and they tend to
have similar activity. Now do a \emph{random} train/test split: for almost every
test molecule, a near-twin sits in the training set. Prediction becomes trivial,
and you wildly \textbf{underestimate the generalisation error} --- the model looks
great in cross-validation and then fails on genuinely new chemistry.

\begin{intbox}
Random splits answer the question ``can the model interpolate within families it
has already seen?'' But the question we actually care about in drug discovery is
``can it generalise to \emph{new} scaffolds?'' Those are very different questions,
and a random split silently swaps the hard one for the easy one.
\end{intbox}

\begin{starbox}[title={Exam-critical: honest splits and the defensible workflow}]
This is the most-tested QSAR idea. \textbf{Compound-series bias:} because datasets
are full of close analogues, a \emph{random} split puts near-twins in both train
and test, so the test set is partly memorised and generalisation error is
\emph{under}estimated --- a random split does \emph{not} ``eliminate bias'', it is
the bias. Keep it distinct from generic feature/hyperparameter leakage and from
class imbalance --- different problems. \textbf{Better splits:} cluster/scaffold
and temporal splits make the test set structurally (or chronologically) distant.
\textbf{Defensible workflow:} \emph{split first}, then do \emph{all} feature
selection and hyperparameter tuning on the training data only, and touch the test
set \emph{once}. Selecting features on the full dataset before splitting, or tuning
on the test set, both leak.
\end{starbox}

\subsection{Smarter Splits and Cross-Validation}

To estimate the performance we actually care about, make the test set
\emph{structurally distant} from the training set:

\begin{itemize}
  \item \textbf{Random split} --- the optimistic baseline above.
  \item \textbf{Cluster split} --- cluster molecules by structural similarity (with
        a chemical distance like Tanimoto, via sphere-exclusion, hierarchical, or
        Butina clustering, or by Bemis--Murcko scaffold), then keep whole clusters
        on one side of the split. This guarantees a minimum train/test distance and
        is far more realistic.
  \item \textbf{Temporal / time-series split} --- order data by the (first) year
        the activity was measured and predict the future from the past. This
        ``mimics'' how a real drug-discovery project unfolds over time.
\end{itemize}

The same logic carries into cross-validation: \textbf{cluster CV} puts entire
clusters of similar molecules into the same fold, and \textbf{temporal CV} uses an
expanding time window (train on everything up to year $t$, test on year $t{+}1$).

\subsection{\starsec{Multi-Task QSAR and Data Leakage}}

When one model predicts activity against \emph{many} targets at once (multi-task
QSAR), a subtler leak appears: the \emph{same} compound can sit in the training set
for one task and in the test set for another. Information bleeds across tasks and
you overestimate performance on new compounds. The lecture frames three split
strategies and what they trade off:

\begin{center}
\small
\begin{tabular}{lccc}
\toprule
\textbf{Split strategy} & \textbf{No leakage} & \textbf{Task balance} & \textbf{Subset dissimilarity} \\
\midrule
Each task split separately            & no  & yes & ? \\
Molecules split across all tasks      & yes & no  & ? \\
\;\;\dots\ with balancing/stratification & yes & yes & ? \\
\bottomrule
\end{tabular}
\end{center}

The priority order (most to least important) is: \emph{no data leakage} first,
then \emph{task balance}, then \emph{subset dissimilarity}. Splitting molecules
(not rows) across all tasks, with stratification to keep tasks balanced, is the
sweet spot.

\section{\starsec{Evaluation and Assessment}}

Once a model is trained you have to grade it --- and with imbalanced, ranking-style
problems the choice of metric is not cosmetic. The standard ML metrics all apply,
but \emph{which} one you trust depends on what you are using the model for: a
ranked list for virtual screening (where \textbf{ranking matters}), removing toxic
compounds, active learning, and so on.

\subsection{Regression Metrics}

For continuous activity (e.g.\ pIC$_{50}$):

\begin{itemize}
  \item \textbf{Accuracy of the value}: MAE (mean absolute error) and RMSE (root
        mean squared error).
  \item \textbf{Interpretability}: $R^2$, the coefficient of determination.
  \item \textbf{Ranking}: Spearman's rank correlation $\rho$ --- the right metric
        when you only care about getting the \emph{order} right, which is exactly
        the VS use case.
\end{itemize}

\subsection{\starsec{Classification Metrics}}

For active/inactive prediction, start from the confusion matrix (TP, FP, FN, TN).
Two families matter:

\textbf{Threshold (``binary'') metrics} --- computed after you commit to a decision
threshold: accuracy, \emph{balanced} accuracy, precision (PPV), recall /
sensitivity (TPR), specificity (TNR), $F_1$, and Matthews' correlation coefficient
(MCC).

\textbf{Ranking (``probabilistic'') metrics} --- computed across all thresholds, so
they grade the model's scoring directly:

\begin{itemize}
  \item \textbf{AUC-ROC} (area under the receiver operating characteristic, ROC,
        curve) --- the model's ability to separate positives from
        negatives across all thresholds.
  \item \textbf{AUC-PR} (area under the precision--recall curve) /
        \textbf{Average Precision} --- the precision/recall
        trade-off, much more informative than ROC under heavy imbalance.
  \item \textbf{BEDROC} (Boltzmann-enhanced discrimination of ROC) --- a variant of
        ROC that rewards \emph{early recognition}
        (getting actives near the top of the list).
  \item \textbf{Enrichment Factor (EF)} --- how much better than random the true-positive
        rate is within the top-$k$ predictions.
\end{itemize}

\newpage

\begin{starbox}[title={Exam-critical: which metrics reward \emph{early} recognition}]
In virtual screening only the very top of the ranking is ever tested, so the
metric must reward getting actives there. \textbf{Only BEDROC and the Enrichment
Factor are early-recognition metrics} --- BEDROC by exponentially weighting the top
of the list, EF by counting actives in the top $x\%$. The trap: \textbf{AUC-ROC},
\textbf{AUC-PR}, \textbf{Average Precision} and the \textbf{precision--recall curve}
all summarise the \emph{whole} ranking and weight every position (roughly) equally,
so none of them is early-recognition-specific --- AUC-ROC is the classic ``odd one
out''. Concretely: to pick a model when only the top $1\%$ can be assayed, use EF
near the $1\%$ cutoff or BEDROC --- \emph{not} accuracy at a $0.5$ threshold and
\emph{not} global AUC-ROC.
\end{starbox}

\begin{notebox}[title={Imbalance breaks accuracy}]
Bioactivity classification is overwhelmingly inactive-heavy, so a lazy model that
calls \emph{everything} negative can score 99\% accuracy while being useless.
Always reach for imbalance-robust metrics --- precision/recall, $F_1$, MCC, AUC-PR
--- and remember BEDROC/EF when what you really want is \emph{good actives at the
top of the ranking}.
\end{notebox}

\section{\starsec{Target Prediction vs.\ Target Identification}}

So far one model predicts activity against \emph{one} target. Real molecules
interact with many proteins, which raises two related but distinct questions.

\textbf{Drug target prediction} asks: \emph{for a molecule, which target(s) does it
bind, and how strongly?} Classical QSAR is single-task (affinity to one target).
\textbf{Multi-task QSAR} solves it with one network that takes a molecule in and
emits one \emph{sigmoid} output per target; because related targets share
structure, the network can borrow statistical strength across them (and it
gracefully handles the missing labels that riddle the activity matrix). Knowledge
graphs are an alternative route.

\textbf{Target identification} (a.k.a.\ \emph{deconvolution} or
\emph{target fishing}) asks the complementary question: \emph{given a molecule and
several possible targets present at once, which one will it most likely bind?}
That is a \emph{multi-class} problem, so the network uses a \emph{softmax} output
over targets.

\begin{intbox}
The output layer encodes the question. \textbf{Sigmoid} units = ``is it active at
\emph{each} of these targets, independently?'' (multi-label / multi-task).
\textbf{Softmax} units = ``which \emph{one} of these competing targets wins?''
(multi-class / target identification). Same molecular encoder, different head,
different scientific question.
\end{intbox}

\begin{starbox}[title={Exam-critical: target \emph{prediction} vs.\ \emph{identification}}]
\textbf{Target identification} (= deconvolution = target fishing): given one
molecule and several candidate proteins present at once, find the \emph{single most
likely} target --- a one-of-$K$ \textbf{multi-class} problem (softmax). \textbf{Target
prediction}: given a molecule, predict activity against \emph{each} target
independently (any number may be ``yes'') --- a \textbf{multi-task / multi-label}
problem (sigmoid). Both traps invert this: ``deconvolution finds \emph{all} binding
proteins'' (no --- that is prediction) and ``\emph{multi-task} does deconvolution''
(no --- multi-class does). Mirror trap: ranking a whole library against \emph{one}
protein is not deconvolution at all --- that is virtual screening.
\end{starbox}

\section{\starsec{Drug Synergy}}

A final, different prediction task: combinations. \textbf{Synergy} means the
combined effect of two (or more) compounds is \emph{larger than the sum} of their
separate effects (super-additive); the mirror case, \textbf{antagonism}, is when
the combined effect is \emph{smaller than the sum} (sub-additive), and a combined
effect that simply \emph{equals} the additive expectation means ``no interaction''.
Predicting whether a pair is \emph{synergistic} or \emph{antagonistic} is hugely
useful --- combination therapy is standard for
cancer and HIV, while some pairings backfire (grapefruit juice antagonises many
drugs and can make them toxic).

How do we even get a label? One common definition is \textbf{Loewe synergy}, which
decomposes the observed effect of a combination as
\[
  \underbrace{E}_{\text{observed}} \;=\;
  \underbrace{A}_{\text{additive (expected)}} \;+\;
  \underbrace{S}_{\text{synergy}} .
\]
You measure single-agent dose--response curves and the combination's response
surface (e.g.\ a $4\times4$ grid of concentrations across many cell lines), fit
the additive expectation, and read off the synergy $S$.

\begin{notebox}[title={Two traps specific to synergy models}]
\textbf{(1) Permutation invariance.} The prediction for the pair $(A,B)$ must equal
the prediction for $(B,A)$ --- the model architecture has to be invariant to the
order of the molecules.\\[2pt]
\textbf{(2) Data leakage, again.} A random split tests performance on combinations
you have \emph{already} partly explored. To estimate performance on genuinely
\emph{novel} combinations you must split by combination --- ``leave drug
combinations out'' (and, more strictly, leave whole drugs or whole cell lines out).
\end{notebox}

\begin{starbox}[title={Exam-critical: synergy vs.\ antagonism, plus symmetry}]
Two facts, often combined in one trick question. \textbf{Definition:} synergy =
combined effect \emph{greater} than additive; antagonism = combined effect
\emph{smaller} than additive (a distractor that calls the super-additive case
``antagonism'', or the sub-additive case ``synergy'', is wrong). \textbf{Property:}
a sensible synergy model must be \textbf{permutation-invariant} --- predicting on
$(A,B)$ and $(B,A)$ must give the same answer. The fully-correct option pairs
\emph{super-additive} with \emph{order-invariant}; watch for options that get the
definition right but make the model order-\emph{dependent}, or vice versa. Loewe
additivity is the standard ``no-interaction'' reference the comparison is made
against.
\end{starbox}

\begin{intbox}
Notice the recurring villain of this chapter: \textbf{data leakage}. It shows up as
random splits over compound series, as compounds shared across multi-task splits,
and as already-seen drug combinations. Each fancy split (cluster, temporal,
leave-combinations-out) is the same fix applied to a new disguise of the same
problem --- making sure the test set asks a question the training set has not
already answered.
\end{intbox}

This is where the core cheminformatics arc closes: we can represent molecules,
screen libraries, build QSAR models, evaluate them honestly, and extend them to
multiple targets and combinations. The \emph{advanced} methods --- message-passing
graph neural networks, multi-modal and contrastive models, active and few-shot
learning, and molecule generation --- build directly on this foundation, and are
the subject of the chapters that follow.


% ==========================================================
%  CHAPTER 3 — ADVANCED METHODS FOR QSAR   (Lecture 03)
% ==========================================================
\chapter{Advanced Methods for QSAR}
\label{chap:advanced}

The QSAR models of Chapter~\ref{chap:vsqsar} treated molecules as fixed feature
vectors (descriptors, fingerprints) and fed them to off-the-shelf learners. This
chapter upgrades both ends: representations that are \emph{learned} from the
molecular graph itself, and learning settings tailored to the small, biased,
multi-modal data that life sciences actually produce --- graph neural networks,
proteochemometrics and contrastive learning, few-shot and active learning, and
multiple-instance learning. Zooming out, deep models for molecular property
prediction fall into a few lineages --- \emph{descriptor-based} (fixed feature
vectors, as in Chapter~\ref{chap:vsqsar}), \emph{graph-based} (representations
learned from the molecular graph, the focus here), and \emph{SMILES/string-based},
with \emph{hybrid} models combining them --- a useful map for the methods that
follow.

\section{\starsec{Molecules as Graphs}}

A molecule is naturally a \textbf{graph} $G=(\mathcal{V},\mathcal{E})$: atoms are
nodes, bonds are edges, and both can carry labels. Graphs are extremely general
--- images (pixels + adjacency), time series, and sets are all special cases ---
and their defining nuisance is that \emph{there is no canonical ordering of the
nodes}. In machine-learning notation a graph is a pair $(\bm{X}, \bm{A})$: a
\textbf{node-feature matrix} $\bm{X}\in\R^{n\times d}$ (one feature row per atom)
and an \textbf{adjacency matrix} $\bm{A}\in\R^{n\times n}$ encoding which atoms are
bonded. Edges (bonds) can carry their own feature matrix too.

\begin{intbox}
Keep that ``$(\bm{X},\bm{A})$'' picture in your head: a molecular graph is
\textbf{two matrices} --- one for the atoms (node features) and one for the bonds
(the edges). Typical node features are the \emph{atomic number}, the
\emph{hybridization state}, the \emph{formal charge}, and so on; bond type
(single/double/aromatic) lives in the edge features. A graph model can and does
use \emph{both} --- it is not limited to node features.
\end{intbox}

Because node order is arbitrary, we want two symmetries: node-level
representations should be \textbf{equivariant} (permute the atoms $\to$ the
per-atom outputs permute the same way), and graph-level predictions should be
\textbf{invariant} (permute the atoms $\to$ the molecule's predicted property is
unchanged). Formally a function $g$ is \emph{invariant} under a group operation
$u$ if $g(u(\xvec)) = g(\xvec)$, and \emph{equivariant} if
$g(u(\xvec)) = u(g(\xvec))$. Permutation-invariant pooling (sum/mean/max) is the
standard way to get graph-level invariance.

\newpage

\begin{starbox}[title={Exam-critical: graphs in, symmetries out}]
A molecular graph is \textbf{two matrices} --- a node-feature matrix $\bm{X}$ whose
atom features are \emph{atomic number, hybridisation, formal charge}, and an
\textbf{adjacency/edge matrix} $\bm{A}$. Bond type (single/double/aromatic) lives in
the \textbf{edge features}, and a GNN uses \emph{both} node \emph{and} edge
information --- ``only node features'' and ``bond type cannot be included'' are both
false. \textbf{Relabel} the atoms without changing connectivity: node
representations are \textbf{equivariant} (they permute the same way), the
graph-level prediction is \textbf{invariant} (unchanged --- the same graph,
reindexed). The trap is ``the prediction may change because $\bm{A}$ was reordered''.
\end{starbox}

\section{\starsec{Graph Neural Networks}}

\subsection{The Message-Passing Framework}

A \textbf{message-passing neural network (MPNN)} learns atom representations by
repeatedly letting each atom ``talk to'' its bonded neighbours. Starting from
$\bm{h}_i^0 = \xvec_i$ (the atom's input features), each round does three steps:

\begin{thmbox}[title={Message passing (one layer)}]
\begin{align*}
\text{message:}\quad   & \bm{m}_{ij}^{t+1} = \phi\!\left(\bm{h}_i^t,\bm{h}_j^t,\bm{e}_{ij}\right),\\
\text{aggregate:}\quad & \bm{m}_i^{t+1} = \textstyle\sum_{j\in N(i)} \bm{m}_{ij}^{t+1},\\
\text{update:}\quad    & \bm{h}_i^{t+1} = \psi\!\left(\bm{h}_i^t,\bm{m}_i^{t+1}\right),
\end{align*}
where $\phi,\psi$ are small MLPs and the aggregation must be
\emph{permutation-invariant} (it sums over an unordered neighbourhood).
\end{thmbox}

\begin{starbox}[title={Exam-critical: the one-layer order is message $\to$ aggregate $\to$ update}]
One MPNN layer is exactly \textbf{message generation $\to$ aggregation $\to$
update}, in that order --- the \textbf{update must come last}. ``Graph convolution''
is \emph{not} a separate sub-step (a GCN is a special case of the whole layer), and
aggregation is \emph{not} the final step: orderings like ``message, graph
convolution, aggregation'' or ``graph convolution, message, update'' are the traps.
\end{starbox}

After $T$ rounds, a permutation-invariant \textbf{read-out}
$\hat{y} = R(\{\bm{h}_1,\dots,\bm{h}_N\})$ pools the atom vectors into one
molecule vector for the final prediction --- e.g.\
$\hat y = \bm{\omega}^{\!\T}\frac{1}{|G|}\sum_{v\in G}\bm{h}_v$. The choice of
aggregation matters: \emph{sum} is the most expressive but propagates signal
poorly, \emph{mean}/\emph{max} are smoother but lose information; variance-preserving
aggregation is a recent compromise.

\newpage

\begin{exbox}[title={Message passing by hand: a 4-atom chain}]
\emph{Intuition first --- think of \textbf{gossip}.} Each atom knows a little; every
round it tells its bonded neighbours what it currently knows, listens to \emph{all}
of them at once, and updates itself. After $T$ rounds each atom has heard from
everything within $T$ bonds.

Take the chain $A\!-\!B\!-\!C\!-\!D$ with one-number features
$h_A^0{=}1,\ h_B^0{=}2,\ h_C^0{=}3,\ h_D^0{=}4$ and neighbourhoods
$N(A){=}\{B\}$, $N(B){=}\{A,C\}$, $N(C){=}\{B,D\}$, $N(D){=}\{C\}$. Use the simplest
possible layer --- \emph{message} $=$ the neighbour's value, \emph{aggregate} $=$
\textbf{sum}, \emph{update} $=$ ``add it to myself'':
$h_i^{t+1}=h_i^t+\sum_{j\in N(i)}h_j^t$. (A real GNN replaces these with small learned
MLPs and a nonlinearity, but the bookkeeping is identical --- and \textbf{what gets
aggregated is exactly the neighbours' current vectors}.)

\textbf{Round 1} (each atom $+$ sum of its neighbours):
\[
h_A^1{=}1{+}2{=}3,\quad h_B^1{=}2{+}(1{+}3){=}6,\quad
h_C^1{=}3{+}(2{+}4){=}9,\quad h_D^1{=}4{+}3{=}7.
\]
Every atom now encodes \emph{itself $+$ its direct neighbours}.

\textbf{Round 2} (reuse the round-1 numbers):
\[
h_A^2{=}3{+}6{=}9,\quad h_B^2{=}6{+}(3{+}9){=}18,\quad
h_C^2{=}9{+}(6{+}7){=}22,\quad h_D^2{=}7{+}9{=}16.
\]
Key point: $h_A^2{=}9$ now depends on $C$, which is \textbf{two bonds away} ---
information travelled $2$ hops in $2$ rounds ($T$ layers $\Rightarrow$ a $T$-hop
receptive field).

\textbf{Read-out:} pool the atom values into one molecule vector --- e.g.\ sum
$=9{+}18{+}22{+}16{=}65$ --- and feed that to an MLP to predict the molecule's
property.
\end{exbox}

It helps to place MPNNs on a spectrum of how a node combines its neighbours:
\emph{convolutional} layers use \textbf{fixed} coefficients (like a CNN; a GCN is
exactly this case), \emph{attentional} layers use \textbf{learned attention
weights} (like a Transformer), and full \emph{message-passing} layers compute a
\textbf{learned message} per edge --- the most general of the three.

A \textbf{graph convolutional network (GCN)} is a restricted MPNN that does
everything in one matrix step,
$\bm{H}^{l+1} = \sigma\!\big(\tilde{\bm{D}}^{-1/2}\tilde{\bm{A}}\tilde{\bm{D}}^{-1/2}\bm{H}^l\bm{W}^l\big)$
with $\tilde{\bm{A}}=\bm{A}+\bm{I}$ (self-loops) --- here all neighbours share the
same weight. The practical QSAR workhorse is \textbf{ChemProp}, a \emph{directed}
MPNN (D-MPNN) that passes messages along bonds (each undirected bond becomes two
directed edges); for best performance it is combined with classical descriptors.
For 3D tasks, \textbf{E($n$)-equivariant} GNNs (EGNN) additionally track atom
coordinates and stay equivariant to rotations/translations. (The equivariance is
cheap to see: a message built from \emph{interatomic distances}
$\norm{\xvec_i-\xvec_j}$ is unchanged by any rotation $\bm{R}$ and translation
$\bm{t}$, because $\norm{\bm{R}\xvec_i+\bm{t}-(\bm{R}\xvec_j+\bm{t})}=\norm{\xvec_i-\xvec_j}$.)

\begin{notebox}[title={A zoo of 3D message-passing layers --- and their symmetry group}]
The same message/aggregate/update skeleton yields different \emph{geometric}
guarantees depending on what the message is allowed to see:
\emph{plain GNN} (uses only node/edge features) --- \textbf{non-equivariant} to
geometry; \emph{Radial Field} and \emph{EGNN} (use distances) ---
\textbf{E($n$)-equivariant}; \emph{SchNet} (uses distances, predicts a scalar) ---
\textbf{E($n$)-invariant}; \emph{TFN} (Tensor Field Networks, uses directional
tensors) --- \textbf{SE($3$)-equivariant}. Rule of thumb: feed the network only
\emph{invariant} geometric quantities (distances, angles) and the output is
automatically invariant; track \emph{coordinates} and update them with
direction vectors and you get equivariance.
\end{notebox}

\subsection{\starsec{Where GNNs Struggle}}

GNNs are not a free lunch, and the lecture is explicit about their failure modes:

\begin{itemize}
  \item \textbf{Limited expressivity} --- message passing is at most as powerful as
        the Weisfeiler--Lehman graph-isomorphism test, so some non-isomorphic
        graphs are indistinguishable.
  \item \textbf{Oversmoothing} --- too many message-passing rounds make all atom
        representations collapse to the same value.
  \item \textbf{Oversquashing} --- the influence of distant atoms decays
        exponentially with path length, so \textbf{long-range interactions} between
        atoms are hard to capture.
  \item \textbf{Under-reaching} --- with too few rounds, information simply doesn't
        travel far enough across a large molecule.
\end{itemize}

\begin{starbox}[title={Exam-critical: tell the four GNN pathologies apart}]
The exam tests the \emph{distinction}, not the names.
\textbf{Limited expressivity}: cannot separate some non-isomorphic graphs (bounded
by the Weisfeiler--Lehman test). \textbf{Oversmoothing}: \emph{too many} layers
$\Rightarrow$ node embeddings converge, so \emph{local} atom environments become
indistinguishable. \textbf{Oversquashing}: distant signal is crushed through narrow
graph \emph{bottlenecks}, so \emph{long-range} interactions are lost.
\textbf{Under-reaching}: \emph{too few} layers $\Rightarrow$ information never
reaches distant atoms. Mind the mirror image --- oversmoothing is the ``too many
layers'' failure, under-reaching the ``too few'' one; oversquashing is about
bottlenecks, not layer count.
\end{starbox}

\begin{notebox}[title={GNN vs.\ Random Forest --- a realistic comparison}]
GNNs are not ``always better''. On the small, feature-rich datasets typical of
QSAR, a \textbf{random forest} often \emph{overfits less} and is \emph{more
interpretable} (you can read off feature importances) than a GNN, while GNNs are
compute-hungry, struggle with very large/complex graphs and long-range effects,
and are \emph{not} less sensitive to noisy data. The honest verdict from
Chapter~\ref{chap:vsqsar} still holds: no single architecture wins everywhere.
\end{notebox}

\section{\starsec{Using Target Information: Proteochemometrics and Multimodal Models}}

Plain QSAR has \textbf{no target awareness}: $f_{\params}(\text{molecule})$ knows
nothing about \emph{which} protein it is predicting activity for, so it cannot
generalise to a new target. \textbf{Proteochemometrics (PCM)} fixes this by giving
the model \emph{both} the molecule and a representation of the target:
$f_{\params}(\text{molecule},\text{target})$. The target can be encoded by
handcrafted amino-acid descriptors, learned from the sequence, or taken from a
pretrained protein language model (e.g.\ ESM). Pooling data across targets gives
\emph{more} training data and lets the model learn cross-target correlations, at
the price of limited applicability to truly novel targets. PCM is one instance of
a \textbf{multimodal model} --- an architecture that integrates several modalities
(molecule, protein, image, text) and fuses them early, jointly, or late to exploit
their complementary information. The point worth absorbing is how \emph{many}
modalities a single molecule offers: \emph{structural} (2D graph, SMILES/InChI
strings, 3D geometry, fingerprints, descriptors, electron density, line drawings),
\emph{textual} (assay descriptions, free scientific text), \emph{reaction-based}
(reactions, reaction templates), and \emph{activity-based} (bioassay-activity
fingerprints) --- alongside accompanying modalities such as the protein target,
microscopy images, and organism. Multimodal models exist precisely to fuse several
of these.

Concretely, PCM architectures have evolved from \emph{concatenating non-learned
molecule and target encodings} at the input of a classical learner (random forest,
SVM, FNN) to fully \emph{learned} fusion: \textbf{DeepPCM} bottlenecks the compound
and protein descriptors separately and fuses them late; \textbf{DeepDTA} and
\textbf{DeepCDA} encode the SMILES string and the amino-acid sequence with
1D-convolutions (the latter adding an LSTM and a two-sided attention mechanism). A
notable recent twist is \textbf{HyperPCM}: a \emph{hypernetwork} reads the target
and \emph{outputs the weights} of a target-specific QSAR model
$f(\xvec;\params)$ --- which, because the target conditions the model directly,
also enables zero-shot prediction on unseen targets (revisited in
Section~\ref{sec:littledata}).

\begin{starbox}[title={Exam-critical: what makes a model proteochemometric}]
PCM \textbf{jointly encodes the molecule \emph{and} a representation of the protein
target}, $f_{\params}(\text{molecule},\text{target})$, so a single model spans many
targets. Two distractors to reject: a \emph{softmax over candidate targets from the
molecule alone} has \textbf{no target representation}, and a \emph{separate
prediction head per target} (plain multi-task) still takes \textbf{molecular inputs
only} --- neither is PCM.
\end{starbox}

\subsection{\starsec{Contrastive Learning}}

A particularly powerful multimodal recipe is \textbf{contrastive learning}, a
\emph{self-supervised} technique that learns by pulling matching pairs together
and pushing mismatched pairs apart in a \textbf{shared latent space}. CLIP is the
canonical example (images $\leftrightarrow$ captions); in drug discovery the pairs
are \emph{(molecule, X)} where $X$ is a binding pocket, a protein sequence, a cell
image, or an assay description (CLAMP, DrugCLIP, CLOOME, SPRINT, ConGLUDe). The
amount of paired data varies enormously by modality --- structure-based
$\sim$10K pairs, imaging-based $\sim$1M, assay-based $>$100M --- and it is exactly
this large scale that lets the learned embeddings support \emph{zero-shot}
generalisation to new tasks.

\begin{notebox}[title={Contrastive learning --- the exact wording}]
Contrastive learning encodes multiple modalities into a \textbf{shared latent
space} and is \textbf{self-supervised} (not unsupervised). It \emph{aligns}
representations of similar/dissimilar pairs; it does not ``concentrate
representations to make predictions'' --- that description is the distractor.
\end{notebox}

\section{Learning from Little Data}
\label{sec:littledata}

\subsection{Few-Shot and Zero-Shot Learning}

A brand-new target may come with $<10$ measured actives --- far too few for a
standard QSAR/PCM model. \textbf{Few-shot learning} explicitly adapts to a small
labelled \emph{support set}, while \textbf{zero-shot} makes predictions with
\emph{no} task-specific data at all. The main routes are
\emph{pretraining + fine-tuning / transfer learning}, \emph{in-context /
retrieval-based} methods that condition on example compounds without weight
updates (e.g.\ MHN-FS, a Modern Hopfield Network), and \emph{meta-learning} (e.g.\
MAML) which trains across many tasks to learn rapid adaptation. Empirically the
picture is humbling: on the FS-Mol benchmark, simple \emph{fine-tuning/probing}
baselines (a linear or quadratic probe on a pretrained encoder such as CLAMP) are
competitive with --- and sometimes beat --- elaborate meta-learning methods, so a
strong baseline is worth trying first.

\subsection{Active Learning}

When labels are expensive, don't label randomly --- let the model pick.
\textbf{Active learning} iteratively trains a model, has it select the
\emph{most informative} molecules from the screening library, sends those for
experimental labelling, and retrains. ``Informative'' usually means \emph{high
prediction uncertainty} (near the decision boundary or in under-explored chemical
space), often combined with \emph{diversity filters} to avoid picking near-identical
compounds. Uncertainty itself is estimated by ensembling, Monte-Carlo dropout, or
evidential learning, and should be \emph{calibrated} (expected calibration error
(ECE), Brier score, reliability
diagrams).

\section{\starsec{Multiple Instance Learning}}

Sometimes the label belongs not to a single molecule but to a \emph{set}.
\textbf{Multiple instance learning (MIL)} trains on \textbf{bags}
$X=\{\xvec_1,\dots,\xvec_M\}$ that carry a \emph{single} label $y$ --- the
individual instances are \emph{not} labelled. The classic chemistry motivation
(Dietterich, 1997) is a molecule with many 3D \emph{conformations}: the molecule is
active if \emph{some} conformation binds, but we don't know which one. The same
framing fits \emph{compound mixtures} (e.g.\ a plant extract is one ``bag'' of many
molecules with a single measured activity).

\begin{notebox}[title={What MIL requires (and what it doesn't)}]
A MIL model must be \textbf{invariant to permutations} of the instances,
$\hat y = g(\{\xvec_1,\dots,\xvec_M\}) = g(\pi(\{\xvec_1,\dots,\xvec_M\}))$, and
must \textbf{handle variable bag sizes}. The bag has \emph{one} label (not one per
instance), and the size is \emph{not} fixed. The fraction of instances that
actually indicate the bag's class is the \emph{witness rate}.
\end{notebox}

Architecturally, MIL uses ``Deep-Sets''-style models: encode each instance, then
apply a \emph{permutation-invariant pooling} (mean/max/sum/attention) and an output
layer. \emph{Bag-level} scoring (pool representations, then predict) usually beats
\emph{instance-level} scoring (predict per instance, then pool). Attention-based MIL
learns which instances matter.

\section{Large Language Models for Molecules}
\label{sec:llms}

The lecture closes with large language models (LLMs) as an emerging tool in the
molecular sciences. The trajectory of (bio)molecular LLMs runs from biomedical-text
models toward multimodal systems that combine text with molecules and proteins. On
knowledge-heavy chemistry benchmarks --- notably \textbf{ChemBench} (thousands of
curated questions across many chemistry subfields) --- leading general-purpose LLMs
(e.g.\ GPT-4, Claude~3) already rival or exceed the \emph{average} human chemist.
The catch is \emph{structure}: on tasks that require genuinely reading a molecule
--- name/format conversion (SMILES $\leftrightarrow$ IUPAC name / molecular formula)
and quantitative property prediction (solubility, lipophilicity, blood-brain-barrier
permeability, toxicity) --- general LLMs fall well behind task-specific, non-LLM
models, and domain fine-tuning (e.g.\ \textbf{LlaSMol}) narrows but does not close
the gap. Bottom line: promising for chemical \emph{reasoning} and knowledge
retrieval, not yet reliable for structure-sensitive prediction.


% ==========================================================
%  CHAPTER 4 — GENERATIVE MODELS FOR MOLECULES   (Lecture 04, 04_Generative_models)  [PART 1 of 2]
%  >>> Paradigms + evaluation. The part-2 material (architectures: LLMs, advanced
%      constraints, GFlowNets, latent-variable models, diffusion/3D) was split out into
%      Chapter 5 (Generative Model Architectures). Mapping is now 1:1: Ch.4=Lec04,
%      Ch.5=Lec05, Ch.6=Lec06.
% ==========================================================
\chapter{Generative Models for Molecules}
\label{chap:generative}

So far every model has been \emph{discriminative}: given a molecule, predict a
number. Generative models run the arrow the other way --- \emph{produce} new
molecules with desired properties. This is \textbf{de novo drug design}: where
virtual screening is $f_{\params}(\text{molecule}) = \text{property}$, generation
is $\text{molecule} = g_{\bm\gamma}(\text{property})$. The motivation is the sheer
size of chemical space: $\sim$$10^{8}$ molecules have ever been synthesised, but
there are perhaps $\sim$$10^{60}$ drug-like ones --- screening can only ever look
at a sliver, so why not \emph{design} directly? Put another way, virtual screening can
only rank molecules someone has already enumerated --- a library of perhaps
$\sim$$10^{6}$ screened compounds --- whereas de novo design tries to reach into the
vast unexplored remainder of that $\sim$$10^{60}$ space.

\section{\starsec{How to Generate a Molecule}}

Generators differ in the \emph{representation} they emit:

\begin{itemize}
  \item \textbf{String-based} --- generate a SMILES or SELFIES string token by
        token (autoregressive RNN/LSTM/transformer).
  \item \textbf{Graph-based} --- build the molecular graph directly, atom by atom
        (autoregressive) or all at once (matrix-based: sample an adjacency tensor
        and node matrix). Must enforce valence and connectivity, so chemical validity
        is hard and the action space can be large --- but it is natural for
        \emph{substructure edits}.
  \item \textbf{Rule-based} --- apply transformation rules to starting compounds.
  \item \textbf{3D} --- generate atom coordinates (diffusion models, equivariant).
\end{itemize}

The most common starting point is \textbf{autoregressive SMILES generation}: encode
the string in one-hot (chemically meaningful tokens), then predict the next token
with $\hat P_{\params}(\xvec) = \prod_t \hat P_{\params}(x_t \mid x_{t-1},\dots,x_1)$.
An LSTM learns the awkward SMILES syntax surprisingly well.

\begin{notebox}[title={Generative models --- the facts the exam loves}]
\textbf{Autoregressive} models \emph{can} generate SMILES; an LSTM reaches up to
$\sim$\textbf{98\% valid} strings, and with \textbf{SELFIES} validity is
\textbf{100\%}. So the claim that SMILES generators produce ``$<$10\% valid'' is
false. Generative models explore \emph{novel} chemical space (not just the
\emph{known} synthesised set), and \textbf{GANs absolutely can} generate molecules
(alongside VAEs, normalizing flows, and diffusion models).
\end{notebox}

\newpage

\section{\starsec{Three Training Paradigms}}

\begin{description}
  \item[Distribution learning.] Learn to \emph{mimic} the property distribution of a
        large library of real molecules (PubChem/ChEMBL) and reproduce its
        ``syntax'' to emit realistic molecules. Self-supervised; used to expand VS
        libraries and as a \emph{starting point} for goal-directed design. Models:
        autoregressive, VAE, GAN, normalizing flow, diffusion.
  \item[Goal-directed learning.] \emph{Optimise} molecules for a \textbf{specific
        objective} (e.g.\ bioactivity) using a \textbf{scoring function} and an
        iterative loop.
  \item[Conditional generation.] Generate molecules that \emph{explicitly satisfy}
        an input condition (property value, scaffold, 3D pocket) by learning a
        shared latent space of molecules and conditions.
\end{description}

\begin{notebox}[title={Distribution learning vs.\ goal-directed}]
A subtle exam distinction: \emph{distribution learning} aims to create molecules
\textbf{similar to the training set}; \emph{goal-directed} generation aims to create
molecules with \textbf{specific desired properties} and can be optimised with
\textbf{reinforcement learning}. ``Replicating existing structures'' is \emph{not}
the point of goal-directed design.
\end{notebox}

\subsection{\starsec{Goal-Directed Optimisation in Practice}}

The optimisation loop can be hill-climbing (iterative distribution learning),
\emph{genetic algorithms}, \emph{reinforcement learning} (define an agent, generate,
score, reward, retrain --- e.g.\ REINVENT), or continuous optimisation in a VAE
latent space. Real drug design is \textbf{multi-objective} (maximise efficacy,
selectivity, QED, synthesizability; minimise toxicity, lipophilicity), so there is
no single best molecule --- trade-offs are handled either by \emph{score aggregation}
--- collapsing the objectives into one number via a weighted arithmetic mean, a
weighted \emph{geometric} mean (which punishes any single near-zero objective), or a
weighted \emph{max} / Chebyshev score (which focuses on the worst objective) --- or by
\emph{Pareto ranking} (keep non-dominated solutions, without committing to weights). A nasty
failure mode is \textbf{mode collapse}: the generator keeps proposing the same or
very similar molecules, whereas we want \emph{diverse} candidates to maximise
downstream success.

\medskip
\noindent\textbf{Where the score comes from.} The objective is rarely a single oracle;
it is assembled from several sources, combined into one set of \emph{objective functions
and constraints}. At the cheap, fast end sit \emph{physicochemical / drug-likeness}
properties (QED, molecular weight, clogP), \emph{structure-based} models (docking a
candidate into a target pocket), and \emph{machine-learning} predictors (QSAR / activity
models); at the slow, expensive, but most trustworthy end sit \emph{experimental
screening} and \emph{human-expert} judgement. The cheap proxies drive the inner
optimisation loop, while the expensive signals are reserved for validating the
survivors --- which is also exactly why proxy \emph{reward hacking} (above) is such a
risk.

The natural way to \emph{drive} an autoregressive generator toward a goal is to read
it as a \textbf{reinforcement-learning} problem --- building a molecule one step at a
time \emph{is} a sequential decision process. The \emph{state} $s$ is the partial
molecule (a SMILES prefix or partial graph); an \emph{action} $a$ is the next
construction step (next token, atom, bond, fragment, or reaction); the \emph{policy}
$\pi(a\mid s)$ \emph{is} the generative model (a distribution over next actions); an
\emph{episode} $\tau$ is one finished molecule; and the \emph{reward} $R(\tau)$ is the
score of that final molecule (activity, QED, synthetic accessibility (SA),
similarity, \dots). Autoregressive
(causal, next-token) models couple to RL especially cleanly because the policy already
factorises over construction steps --- but the reward structure makes the problem nasty.

\begin{starbox}[title={\ensuremath{\bigstar}\, Why molecular RL is hard}]
\begin{itemize}
  \item \textbf{Delayed reward} --- the score is known only \emph{after} a full
        molecule has been generated, not step by step.
  \item \textbf{Sparse / weak signal} --- early in training most molecules are
        invalid, inactive, or low-scoring.
  \item \textbf{Proxy objective \& reward hacking} --- the reward is a QSAR, docking,
        or heuristic score (\emph{not} experimental success), so the generator can
        exploit \emph{artifacts} of the scoring function.
  \item \textbf{Diversity is not automatic} --- unless explicitly rewarded, RL does
        not try to cover multiple chemical solutions.
  \item \textbf{Distribution shift} --- aggressive optimisation drifts the agent away
        from the realistic chemistry the prior had learned.
\end{itemize}
\end{starbox}

\medskip
\noindent\textbf{The practical RL toolbox (REINVENT-style).} Real systems tame these
problems with a standard kit: \emph{reward shaping} (turn raw scores into a usable
learning signal and balance objectives), \emph{prior anchoring} (a KL-style penalty
keeping the agent close to the pretrained distribution, which curbs unrealistic
chemistry and reward hacking), \emph{chemical filters} (penalise reactive, unstable, or
otherwise undesirable structures), \emph{diversity filters} (penalise repeated
scaffolds / already-explored regions), and \emph{replay memory} (store high-scoring
molecules and reuse them, improving sample efficiency). Two extensions explicitly
\emph{strengthen the learning signal}: \textbf{Augmented Hill-Climb} updates only on the
top-scoring molecules in each batch --- more optimisation pressure and better sample
efficiency, at the \emph{risk} of reduced exploration and faster mode collapse --- and
\textbf{Augmented Memory} pairs replay with SMILES augmentation, showing each
high-scoring molecule to the model through several randomised SMILES to extract more
signal when scoring is expensive.

\medskip
\noindent\textbf{Promoting diversity / beating mode collapse.} Because diversity is
not automatic, a range of practical fixes have been proposed: \emph{select diverse
high-scoring molecules} at each step while gradually raising a property threshold
(Rupakheti et al.); \emph{mix exploitation with exploration} by blending a fixed
exploratory generator with the trainable one during SMILES generation (Liu et al.\ ---
\textbf{DrugEx}, which trains paired \emph{exploitation} and \emph{exploration}
networks and samples from the explorer with probability $\varepsilon$, adding
Pareto-based multi-objective selection for diversity);
\emph{penalise redundant solutions} by docking the reward of molecules that are
similar to past hits kept in a \emph{memory buffer} (Blaschke et al., the diversity
filter in REINVENT); \emph{expand around key substructures} by generating many
completions of extracted scaffolds (Chen et al.); \emph{sample proportionally to
reward} so that multiple modes are covered (Bengio et al.\ --- the GFlowNet idea
below); and \emph{reuse diverse SMILES variants} by training on non-canonical SMILES
of the same molecule (Bjerrum et al.). The common thread is to \emph{explicitly}
reward or inject diversity --- simply taking the single highest-scoring molecule each
round is pure exploitation and does the opposite.

\begin{starbox}[title={\ensuremath{\bigstar}\, Combating analogue rediscovery}]
When a goal-directed generator keeps rediscovering close analogues of earlier hits,
\textbf{genuine remedies all add or reward diversity}: \emph{mixing} a fixed
exploratory generator with the trainable one, or \emph{penalising} reward for
molecules similar to a memory buffer of past high-scorers. The exam \textbf{trap}:
\emph{always selecting the single highest-scoring molecule} as the next training
example is \emph{increased exploitation} --- it \emph{worsens} rediscovery and is
\emph{not} a remedy.
\end{starbox}

\subsection{\starsec{Conditional Generation}}

Conditional generators (conditional VAE/RNN/transformer/diffusion) give
\emph{direct control} over properties without a scoring function and enable fast
one-pass generation, but they need \emph{paired} data (molecules with known
properties), may only \emph{approximately} satisfy the target, generalise poorly to
rare conditions, and risk \emph{property entanglement}.
\vspace{-0.2 em}
\section{\starsec{Evaluating Generated Molecules}}

There is no single score. Basic \emph{diagnostic} checks are
\textbf{validity} (\% syntactically correct), \textbf{novelty} (\% not in the
training set), and \textbf{uniqueness} (\% distinct among generated). Beyond those:

\begin{itemize}
  \item \textbf{Diversity} --- internal diversity (mean pairwise $1-$Tanimoto)
        captures spread but not coverage; \emph{\#Circles} counts how many mutually
        dissimilar molecules fit, capturing chemical-space coverage better.
  \item \textbf{Similarity to a reference set} --- compare property distributions via
        Kullback--Leibler divergence or Kolmogorov--Smirnov distance, or compare in a
        learned ``biologically informed'' space via the Fréchet ChemNet Distance.
\end{itemize}

\begin{thmbox}[title={Fréchet ChemNet Distance (FCD)}]
Map real and generated molecules to Gaussians using the last hidden layer of a
pretrained ``ChemNet'' (a target-prediction LSTM over SMILES); with means
$\bm{m},\bm{m}_w$ and covariances $\bm{C},\bm{C}_w$,
\[
  \mathrm{FCD}^2 = \norm{\bm{m}-\bm{m}_w}_2^2
  + \tr\!\Big(\bm{C}+\bm{C}_w - 2(\bm{C}\,\bm{C}_w)^{1/2}\Big).
\]
So FCD compares the \textbf{means and covariances of learned latent
representations} and measures the distance between the \emph{distributions} of
generated vs.\ real molecules --- \emph{not} distances between individual atoms, and
\emph{not} computed on raw structures. It conveniently picks up many biases at once
(drug-likeness, $\log P$, synthesizability).
\end{thmbox}

Rather than re-implement these metrics ad hoc, the field has converged on a handful of
standardised \textbf{benchmarks} that bundle them together. The canonical ones are
\textbf{GuacaMol} (distribution learning plus single- and multi-objective goal-directed
optimisation tasks), \textbf{MOSES} (distribution learning, molecular diversity, and
avoidance of forbidden substructures), and \textbf{MolOpt} (sample efficiency over $23$
oracle scoring functions) --- with later suites such as \emph{MolScore} (a re-implementation
that aggregates GuacaMol/MOSES/MolOpt and adds biological-activity and synthetic-accessibility
evaluations) and 3D-specific ones like \emph{GenBench3D} (conformer quality against
Cambridge-Structural-Database geometries). The takeaway is that each benchmark fixes a
particular \emph{evaluation axis}: which metric matters depends on whether you care about
distribution matching, optimisation efficiency, or 3D geometric realism.

\medskip
\noindent\textbf{Valid is not enough: desirable and suitable.} A diagnostic check tells you a
molecule is \emph{valid}; goal-directed learning and conditional generation aim higher, for
molecules that are \emph{desirable} (they satisfy \emph{all} the specified objectives, not just
syntactic correctness). But a desirable molecule must also be a \emph{suitable} drug-discovery
candidate: it needs reasonable ADME behaviour (Chapter~\ref{chap:drugdiscovery}) and it has to
be \emph{synthesisable} (Chapter~\ref{chap:synthesis}). Those last two checks are exactly why
generation does not end at the generator. A sobering reality check underlines this: of the
40-plus multi-objective generative methods published between 2020 and 2022, only \emph{one}
was experimentally validated --- in-silico desirability is necessary but far from sufficient.
The next chapter first opens up the
\emph{architectures} that implement these paradigms; suitability and synthesis are then
taken up in Chapter~\ref{chap:synthesis}.


% ==========================================================
%  CHAPTER 5 — GENERATIVE MODEL ARCHITECTURES   (Lecture 05, 05_Generative_models_2_ocred.pdf)
%  >>> Part-2 of the generative-models topic, split out of the old monolithic Ch.4.
%      Sections: Chemistry-Specific LLMs, Advanced Generative Constraints, GFlowNets,
%      Latent-Variable Models & Latent-Space Optimisation (star), Diffusion & 3D Generation (star).
% ==========================================================
\chapter{Generative Model Architectures}
\label{chap:generative2}

Chapter~\ref{chap:generative} framed generation by \emph{paradigm} (distribution
learning, goal-directed, conditional) and by how to \emph{evaluate} it. This chapter
surveys the model \emph{families} that actually implement those paradigms --- the
\emph{representation} a generator emits largely \emph{defines the generation problem}.
The lecture's organising taxonomy pairs each representation with the model families
that emit it:
\begin{itemize}
  \item \textbf{String} (SMILES/SELFIES) $\to$ token sequence: RNNs, Transformers,
        VAEs, GANs.
  \item \textbf{2D graph} $\to$ atoms and bonds: graph VAEs, graph flows, graph
        diffusion, GFlowNets.
  \item \textbf{3D graph} $\to$ atoms and coordinates: equivariant diffusion, flows,
        autoregressive 3D models.
  \item \textbf{Reaction route} $\to$ building blocks and reactions: handled by the
        reaction-based design and synthesis-planning material
        (Chapter~\ref{chap:synthesis}).
\end{itemize}
Three big buckets organise the rest: \textbf{sequential} (autoregressive string/graph
generators, including GFlowNets), \textbf{latent-space} models
(learn a continuous space, sample and decode), and \textbf{denoising} models
(diffusion, especially natural for 3D). We also collect the \emph{constraint tasks}
that all of these can be pointed at.

\section{Sequential (Autoregressive) String Generators}

The oldest and still most common way to emit a molecule is to emit a
\emph{string}, one token at a time. A SMILES (or SELFIES) string is consumed
exactly like a sentence, which is why the whole language-modelling toolbox
transfers directly: the earliest molecular generators were \textbf{RNNs} and
\textbf{LSTMs} trained to predict the next SMILES character, and they were later
joined by \textbf{Transformers}, \textbf{VAEs} and \textbf{GANs} over the same
token sequence. \textbf{MolGPT} (Bagal et al., 2022) is the canonical modern
example: a \emph{conditional Transformer} that generates SMILES conditioned on a
scaffold plus physicochemical properties, in the same spirit as the earlier
conditional RNN (\textbf{cRNN}, Kotsias et al., 2020) conditioned on
fingerprints and QSAR properties.

\begin{notebox}[title={``Chemistry LLM'' does not mean ``understands chemistry''}]
It is tempting to conclude that because a Transformer generates valid SMILES, it
has learned chemical structure. The lecture's own verdict points the other way:
general and chemistry-tuned LLMs are strong on chemical \emph{knowledge and
reasoning} yet measurably weak at \emph{reading a structure} (see
Section~\ref{sec:llms} in Chapter~\ref{chap:advanced}). Generating a well-formed
string and understanding the molecule it denotes are different skills --- do not
let the fluency of the output stand in for structural competence.
\end{notebox}

\newpage

\section{Advanced Generative Constraints}

``Molecular generation'' is \emph{not one task}: the same generators are pointed at
several distinct problems, which differ in \emph{what is held fixed} and \emph{what is
generated}. The constraint defines the use case:
\begin{itemize}
  \item \textbf{Full de novo design} --- generate the whole molecule with minimal
        structural constraints ($\varnothing \to$ molecule). Best for \emph{broad
        exploration} of chemical space when no clear starting series exists.
  \item \textbf{Scaffold decoration} --- keep the central core fixed and generate the
        peripheral substituents / R-groups. This is close to medicinal-chemistry
        \emph{lead optimisation} around a known active scaffold.
  \item \textbf{Linker design (fragment linking)} --- keep two disconnected fragments
        fixed (e.g.\ within a binding pocket) and generate a chemically valid linker
        that bridges them. Useful for \emph{fragment linking}, \textbf{PROTAC} design,
        and preserving key binding motifs.
  \item \textbf{Scaffold hopping} --- the inverse of decoration: keep the peripheral,
        activity-relevant features fixed and \emph{replace the central scaffold} with a
        structurally distinct core. Useful for novelty, \textbf{IP}, selectivity, and
        ADMET optimisation.
  \item \textbf{Reaction-based design} --- generate molecules through available
        \emph{building blocks} and feasible \emph{reaction steps} (amide coupling,
        Suzuki coupling, esterification, Diels--Alder, \dots). This constrains
        generation toward molecules that are more likely \emph{synthesisable} --- a
        direct hand-off to the synthesis-planning ideas of
        Chapter~\ref{chap:synthesis}.
\end{itemize}
Concrete RL realisations exist for several of these: \textbf{LibINVENT} performs
reaction-based scaffold decoration and \textbf{Link-INVENT} does linker design, both by
conditioning an encoder--decoder RNN on a SMILES specification and generating the
missing part, which is then assembled into the full molecule.

\vspace{-0.5 cm}

\section{Generative Flow Networks (GFlowNets)}
Standard reinforcement learning for molecule generation often suffers from \textbf{mode collapse} --- the agent finds a single high-scoring chemical pathway (often by exploiting a loophole in the proxy reward) and repeatedly generates the same or very similar molecules, abandoning chemical diversity.

\textbf{Generative Flow Networks (GFlowNets)} offer a modern alternative. Instead of greedily maximizing the expected reward, GFlowNets treat the sequential generation of a molecule as a network flow. The agent is trained to sample terminal states (completed molecules) with a probability strictly \emph{proportional} to their objective reward, $P(x) \propto R(x)$. This ensures that the model actively explores and samples a diverse set of high-reward \emph{modes} across the chemical space, naturally mitigating mode collapse while still requiring a scoring function to guide the flow probabilities.
\vspace{-0.5 cm}
\section{\starsec{Latent-Variable Models and Latent-Space Optimisation}}

Autoregressive models build a molecule step by step; the \emph{latent-variable}
alternative is to first learn a \textbf{continuous space} of molecules (typically with
a VAE --- but GANs and normalizing flows learn such spaces too) and then generate by
sampling a latent vector $\zvec$ and \emph{decoding} it.
New molecules come from sampling $\zvec$; goal-directed design becomes \emph{moving}
$\zvec$ toward desired properties --- optimisation in a smooth space instead of
surgery on discrete graphs.
The three classic continuous-space families differ in \emph{how} they tie molecules to
$\zvec$: a \textbf{VAE} pairs an encoder $q(\zvec\mid x)$ with a decoder
$p(x\mid\zvec)$; a \textbf{GAN} trains a generator $G(\zvec)$ against a discriminator,
with no explicit encoder; and a \textbf{normalizing flow} uses an \emph{invertible} map
so that $\zvec=f^{-1}(x)$ and $x=f(\zvec)$ exactly.

\textbf{The decoder depends on the representation.} A latent $\zvec$ is turned back
into a molecule by a \emph{sequential SMILES decoder} (token sequence), a \emph{graph
decoder} (atoms + bonds), or a \emph{fragment decoder} (motifs / scaffold pieces).
Graph decoding in particular comes in three flavours:
\begin{description}
  \item[One-shot] --- predict all atom types and the bond-adjacency tensor at once
        (one decoder pass from $\zvec$ to the whole molecule). Conceptually simple, but
        hard to enforce valence and connectivity.
  \item[Sequential] --- $\zvec$ + partial graph $\to$ next atom/bond action; builds
        step by step and can \emph{restrict actions to chemically valid choices} (like
        graph autoregression, but conditioned on $\zvec$).
  \item[Fragment-level] --- $\zvec \to$ chemically meaningful fragments that are
        assembled into the graph; often improves validity and connects naturally to
        scaffold/linker tasks, but is limited by the fragment vocabulary.
\end{description}

\textbf{Why optimise in latent space?} Molecules are \emph{discrete}, but $\zvec$ is
\emph{continuous}. We train (or reuse) a property model $f(\zvec)$, search
$\zvec^\ast = \argmax_{\zvec} f(\zvec)$, then decode $x^\ast \sim p(x \mid \zvec^\ast)$.
Continuity unlocks tools that discrete graphs do not offer: \emph{gradient ascent}
(if $f$ is differentiable), \emph{Bayesian optimisation} (when evaluations are
expensive), \emph{evolutionary search} (mutate/select latent vectors), and
\emph{interpolation} between known actives. The learned prior keeps the search near
realistic chemistry, similar molecules tend to sit close together when the space is
well learned, and sampling many $\zvec$ naturally yields multiple candidates --- at
the cost of only \emph{approximate} control over the decoded structure.

\noindent\textbf{Failure modes of latent-space optimisation.} The smoothness that makes
this attractive is also where it breaks. The latent space may simply \emph{not be
chemically smooth}, so small moves in $\zvec$ can cause large structural changes after
decoding, and high-scoring $\zvec$ can decode to \emph{invalid} molecules. Optimisation
can push $\zvec$ \emph{outside the region seen during training} (a distribution shift
echoing the RL case), and --- exactly as with goal-directed RL --- the
\emph{property predictor $f$ can be exploited} (reward hacking in latent space). Even a
\emph{valid} decoded molecule may be unrealistic or hard to synthesise, which is why
latent optima are filtered and checked rather than trusted directly.

\vspace{-0.5 cm}

\section{\starsec{Diffusion Models and 3D Generation}}

\textbf{Diffusion models generate by denoising}: start from a noisy representation and
learn to \emph{reverse} a gradual noising process, producing a molecule through
\emph{iterative} refinement rather than one-shot decoding. This iterative view is
especially natural for \emph{continuous} representations such as 3D coordinates.

\begin{starbox}[title={\ensuremath{\bigstar}\, Three generation modes, contrasted}]
The exam likes to separate \emph{how} the molecule is produced:
\begin{itemize}
  \item \textbf{One-shot latent decoding} --- map a \emph{single} latent vector
        directly to a \emph{complete} molecule in \emph{one} decoder pass.
  \item \textbf{Autoregressive} --- build the molecule \emph{step by step}, predicting
        the next token (or graph action) from the partial molecule.
  \item \textbf{Diffusion} --- start from \emph{noise} and \emph{iteratively} apply a
        learned \emph{denoising} process toward a molecule.
\end{itemize}
So what distinguishes diffusion from one-shot decoding is the \textbf{iterative
denoising from noise}, \emph{not} a single direct latent$\to$molecule pass (one-shot)
and \emph{not} next-step prediction from a partial molecule (autoregressive).
\end{starbox}

\textbf{What gets denoised matters.} Two regimes:
\begin{description}
  \item[2D graph diffusion] --- generates molecular \emph{identity} (atoms + bonds);
        noise is applied to \emph{discrete} atom/bond types and denoising predicts the
        clean graph. The hard part is \emph{chemical validity}: valence, connectivity,
        aromaticity.
  \item[3D diffusion] --- generates identity \emph{and/or} geometry; noise is applied
        to \emph{continuous} coordinates (sometimes also atom types). The hard part is
        \emph{realistic geometry}: stable structures and correct symmetries, which
        demands rotation/translation-aware models.
\end{description}

\begin{starbox}[title={\ensuremath{\bigstar}\, 3D generation needs equivariance}]
Molecular coordinates have \textbf{no fixed global orientation}, so a 3D generator
must not depend on an arbitrary coordinate frame. If you \emph{rotate or translate}
the input structure, the prediction should rotate/translate \emph{consistently} ---
this property is \textbf{equivariance}, and \emph{equivariant networks} build it in by
construction. Ignore it and the model wastes capacity memorising orientations and
tends to produce unstable, unrealistic geometries.
\end{starbox}

The canonical equivariant backbone is the \textbf{E($n$)-equivariant graph neural
network (EGNN)} (Satorras et al., 2021), which interleaves scalar \emph{node
embeddings} with \emph{coordinate embeddings} and passes messages built from
\emph{pairwise distances} $\norm{\xvec_i-\xvec_j}^2$ (themselves rotation/translation
invariant). Coordinates are then updated only along \emph{relative offsets},
\[
  \xvec_i^{\,t+1} = \xvec_i^{\,t} + C\!\!\sum_{j\neq i}(\xvec_i^{\,t}-\xvec_j^{\,t})\,
  \phi_1(\boldsymbol{m}_{ij}),
\]
which makes each layer \emph{equivariant} to rotations and translations, while the
scalar node features stay \emph{invariant} --- so the whole network respects 3D
symmetry by construction.

\textbf{Conditional 3D diffusion} fixes part of the molecule and denoises the rest:
\begin{itemize}
  \item \textbf{Conformer generation} (e.g.\ GeoDiff): the molecular \emph{identity is
        fixed}; the model denoises \emph{coordinates conditioned on the 2D graph},
        sampling multiple 3D conformers of one molecule. Useful for docking
        preparation, shape-based screening, and 3D-QSAR --- but it is \emph{not} de
        novo molecule generation.
  \item \textbf{Pocket-conditioned ligand generation} (e.g.\ Schneuing et al.): the
        protein \emph{pocket} supplies the 3D context, and the model generates or
        denoises a ligand \emph{inside the binding site}, conditioned on pocket atoms,
        geometry, pharmacophores, or interaction features. It usually needs
        post-processing to recover valid molecules, followed by docking, filtering,
        property prediction, and synthesis checks.
\end{itemize}

\noindent\emph{Two distractors to reject here.} Pocket-conditioned generation
\textbf{conditions the generator on the pocket \emph{during} generation}: the binding
site is an \emph{input} that shapes what gets built. It is \emph{not} ``generate a
ligand with no protein context and only dock afterwards'' --- in that picture the
docking score never steers the generator. And it is \emph{not} conformer generation,
which \emph{fixes the ligand identity} and samples 3D shapes \emph{without} using any
protein information. The pocket drives generation; it is not a post-hoc filter.


% ==========================================================
%  CHAPTER 6 — CHEMICAL REACTIONS AND SYNTHESIS PLANNING   (Lecture 06, 06_Synthesis_ocred.pdf)
%  >>> Renders as Chapter 6 (Ch.5 = Generative Model Architectures now sits above).
% ==========================================================
\chapter{Chemical Reactions and Synthesis Planning}
\label{chap:synthesis}

A generator can dream up a molecule, but a molecule is only useful if someone can
actually \emph{make} it. \textbf{Synthesis planning} --- finding a feasible sequence
of reactions from purchasable starting materials to the target --- is therefore not
a post-processing afterthought but a \emph{core problem} of drug discovery, and it
applies equally to human-designed and AI-designed molecules.

\section{Chemical Reactions}

A \textbf{chemical reaction} transforms one molecule into another by breaking and
forming bonds (electrons move). It is written $a\mathrm{A} + b\mathrm{B} \to
c\mathrm{C} + d\mathrm{D}$ with reactants on the left, products on the right, and
$(a,b,c,d)$ the smallest integers that balance the atom counts (conservation of
mass). Common types are \emph{combination} ($A+B\to AB$), \emph{decomposition}
($AB\to A+B$), \emph{substitution} ($A+BC\to AC+B$), and \emph{metathesis}
($AB+CD\to AD+CB$). Which product actually forms is governed by three factors.
\emph{Thermodynamics} asks whether the transformation is favourable: it depends on
the relative stability / free energy of reactants and products and determines which
products are energetically favoured. \emph{Kinetics} asks whether it happens fast
enough: it depends on the reaction barrier and on conditions such as temperature,
concentration, catalyst, and solvent --- so the \emph{same} reactants under different
conditions can give different products. \emph{Selectivity} asks which product is
actually obtained when several are possible: conditions shift the major product,
yield, and side products. In practice, then, synthesis planning cares not just about
\emph{what} can react but under \emph{which conditions} and with \emph{what
yield/selectivity}.

\textbf{Chemical synthesis} is the artificial execution of useful reactions to obtain
the desired product(s): selecting reactants, usually over multiple steps, from
available or easily accessible building blocks, while choosing and optimising the
reaction conditions.

\section{\starsec{Computer-Assisted Synthesis Planning (CASP)}}

CASP \emph{predates} modern molecular generation: chemists have always needed a
route after proposing a target, and symbolic systems date back to Corey \& Wipke
(1969) and LHASA. Modern ML swapped hand-coded rules for \emph{learned} reaction
models, route scoring, and search. Reactions are stored in a SMILES-like language
with an \textbf{atom-to-atom mapping} (numbering atoms across reactants and
products). Datasets include \textbf{USPTO} (reactions mined from US patents), CAS /
Reaxys / Pistachio, and the Open Reaction Database.

A modern CASP system is a pipeline. From the reaction datasets a \emph{knowledge
base} is curated --- a template library, a reaction network, and curated reaction
sequences --- which feeds two learned components: \emph{reaction prediction} (forward
products, and optionally conditions and selectivity) and \emph{retrosynthesis}. The
planner then recursively traverses a search tree of disconnections, scoring candidate
\emph{routes} and checking each leaf against a \emph{stock database} of purchasable
compounds; a route is complete once every leaf is in stock.

\begin{notebox}[title={Reaction data is biased}]
Most reaction datasets contain only \emph{successful} reactions (no failed
attempts), records are often \emph{incomplete} (missing conditions, yields), and a
recorded route is \emph{not the only ground truth} --- other valid routes exist, so
benchmark labels underestimate good alternative predictions.
\end{notebox}

\section{\starsec{Predicting Reactions}}

Two directions, both supervised ML over reaction data:

\begin{itemize}
  \item \textbf{Forward reaction prediction}: input reactants (and sometimes
        conditions), output the major product.
  \item \textbf{Single-step retrosynthesis} (backward): input a target product,
        output plausible precursors.
\end{itemize}

The two directions are not equally hard. Forward prediction is the most constrained
(reactants $\to$ a single major product) and is generally the \emph{easiest};
single-step retrosynthesis is harder because it is one-to-many (a product has many
valid precursor sets); and \emph{multi-step} planning (next section), which chains
many such decisions under search, is the hardest of the three. Beyond products and
precursors, the same models are also pointed at \emph{other tasks} --- predicting
reaction \emph{conditions} (solvent, catalyst, temperature) and \emph{yield}.

There are two philosophies for how to do it:

\begin{description}
  \item[Template-free] --- treat it as \emph{translation}: a sequence-to-sequence
        ``Molecular Transformer'' directly maps reactant strings to product strings.
        Highly accurate ($\sim$90\%) but hard for a chemist to inspect (no explicit
        rules).
  \item[Template-based] --- ``neural-symbolic AI'': a \emph{neural} part selects a
        \textbf{reaction template}, and a \emph{symbolic} part executes the
        transformation rule encoded in that template.
\end{description}

A \textbf{reaction template} (reaction rule) is a generalised blueprint of a
reaction --- which substructures must be present and how they transform --- designed
by experts or extracted from a database. Templates are extracted from atom-mapped
reactions by recording the bonds that break and form plus a chosen radius of
surrounding atomic environment: a large radius gives a specific template that may not
generalise to new reactants, a small one gives a general template that applies to too
many molecules --- the core complexity-vs-coverage trade-off.

Across these philosophies the lecture distinguishes four ML approaches: rule-based
\emph{expert systems}; \emph{translation} (direct sequence-to-sequence or
graph-to-graph prediction --- the template-free route); \emph{classification}
(predict the reaction template/class with a neural network); and \emph{retrieval}
(learn to fetch the correct template from a database). The last two are both
template-based but differ in how the template is \emph{scored}. Classification uses
$\hat{y} = \softmax(\mathbf{W}\,\hvec^{m}(m))$, where each row of $\mathbf{W}$ is an
independent, freely-adjustable reaction class with no notion of similarity between
templates --- so it degrades on \emph{rare} templates that have few training
examples. Retrieval (MHNreact, built on \textbf{Modern Hopfield Networks}) replaces
$\mathbf{W}$ with learned template embeddings,
$\hat{y} = \softmax(\hvec^{t}(\mathbf{T})\,\hvec^{m}(m))$: reactant and templates are
encoded into a shared associative memory and the closest template is retrieved, so
related templates share structure. Template-free translation (the Molecular
Transformer) is the most accurate ($\sim$90\%) but acts as a \emph{black box} ---
there is no explicit transformation rule or electron-pushing mechanism to inspect.
On the standard \textbf{USPTO-50k} single-step-retrosynthesis benchmark the best
template-based (MHNreact, GLN) and template-free (Molecular Transformer, MEGAN)
models are broadly competitive, with top-1 accuracies around 50\%.

\begin{notebox}[title={AI for chemical reactions --- key claims}]
Some methods rely on \textbf{reaction templates} (transformation rules); deep models
can be \emph{trained on databases of known reactions}; and \textbf{transformer}
architectures are used for reaction prediction. What you do \emph{not} need is
quantum-mechanical data --- ML predicts reactions from reaction databases, not from
QM calculations. And synthesis planning is \emph{not} ``trivial'' nor merely
``selecting reactants'': it designs a whole \emph{sequence} of reactions and uses
\emph{retrosynthetic analysis} to break the target into simpler building blocks.
\end{notebox}
\begin{starbox}[title={Reaction prediction --- exam essentials}]
\textbf{Directions.} Forward prediction maps \emph{reactants $\to$ major product};
single-step retrosynthesis maps \emph{product $\to$ precursors}. Forward is the
\emph{easiest} task, multi-step planning the \emph{hardest}.\\
\textbf{Philosophies.} \emph{Template-free} = sequence-to-sequence translation (the
Molecular Transformer), $\sim$90\% accurate but uninspectable. \emph{Template-based}
= neural-symbolic: a neural part picks a \textbf{reaction template} (a generalised
symbolic rule: required substructures + how they transform), a symbolic part executes
it --- via \emph{classification} $\softmax(\mathbf{W}\hvec^{m})$ or \emph{retrieval}
$\softmax(\hvec^{t}(\mathbf{T})\hvec^{m})$ (MHNreact / Modern Hopfield).\\
\textbf{Data.} Trained on reaction \emph{databases} (USPTO, Reaxys, ORD), \emph{not}
on quantum-mechanical data.
\end{starbox}
\section{\starsec{Multi-Step Retrosynthesis as a Game}}

A full synthesis is many steps, so multi-step retrosynthesis repeats single-step
decisions until every leaf is a purchasable building block. The elegant framing is
as a \textbf{game} played with \emph{Monte-Carlo Tree Search + deep reinforcement
learning} (the AlphaGo recipe):

\begin{thmbox}[title={The retrosynthesis ``game''}]
\begin{itemize}
  \item \textbf{State}: the current set of molecules still to be made.
  \item \textbf{Action}: selecting reactants (i.e.\ choosing a reverse reaction /
        disconnection).
  \item \textbf{Environment}: performs the reaction and returns the new molecules.
  \item \textbf{Reward}: obtained when the current set of molecules are \emph{all}
        available building blocks --- the game is won.
\end{itemize}
\end{thmbox}

\begin{notebox}[title={State vs.\ building blocks --- don't swap them}]
The \emph{state} is the current set of molecules you still need to synthesise; the
building blocks are the \emph{goal} you are trying to reach. Saying ``the state is
the set of available building blocks'' is the classic trap --- it is false.
\end{notebox}

Each move in this game is one round of \textbf{Monte-Carlo Tree Search}, which cycles
through four phases: \emph{selection} (descend the tree to the most promising
unexpanded node via a tree policy), \emph{expansion} (add child nodes by retro-analysing
that molecule into candidate precursor sets), \emph{simulation / rollout} (play out a
route from the new node and evaluate it), and \emph{backpropagation / update} (push the
result back up the path so future selections are better informed).

Because the search tree is astronomically large (like chess or Go --- $\sim$10$^{46}$
and $\sim$10$^{170}$ states respectively), DNNs predict, for unvisited states, both the
\emph{value} (how likely to lead to a solution) and a \emph{policy} (which disconnection
to try). In practice $\sim$50\% of molecules can
be ``solved'' with routes that make sense to chemists --- a ``chemical Turing test''
found expert chemists could not reliably tell AI-planned routes from human ones ---
though whether the routes actually run, and under which conditions, remains uncertain.

This closes the chapter's two halves: \emph{chemical reactions} are predicted with
\emph{supervised} learning over reaction databases, whereas multi-step \emph{chemical
synthesis} is solved as a \emph{reinforcement-learning} search problem.








% ==========================================================
%  CHAPTER 7 — MEDICAL DATA AND IMAGING
% ==========================================================
\chapter{Medical Data and Imaging}
\label{chap:medical}

The second half of the course shifts from molecules to \emph{patients}. The data
look different --- electronic health records, tabular measurements, medical images
--- but the same machine-learning discipline applies, with a few twists that are
specific to clinical data and that the exam likes to probe.

\section{Tabular Medical Data}

A patient table mixes \textbf{data types}, and getting them right drives every
preprocessing choice:

\begin{itemize}
  \item \textbf{Categorical (nominal)} --- unordered labels, e.g.\ smoking status
        (non-/former/current smoker). It is \emph{not} continuous, and you cannot
        take its mean.
  \item \textbf{Ordinal} --- ordered categories without a meaningful numeric
        distance.
  \item \textbf{Continuous (numeric)} --- real-valued measurements like a cholesterol
        level in mg/dL. Being orderable does \emph{not} make it merely ``ordinal''.
\end{itemize}

\begin{notebox}[title={Imputing missing values}]
For a skewed \emph{numeric} column (e.g.\ cholesterol), impute missing values with
the \textbf{median} (robust to skew), not the mean. For a \emph{categorical} column
(e.g.\ smoking status) you cannot use the mean at all --- use the \emph{mode}. Match
the imputation to the data type.
\end{notebox}

\section{Batch Effects}

When samples are processed in different batches (machines, days, labs, reagents),
\textbf{technical} variation gets baked into the data. \emph{Batch effects} create
\textbf{artificial differences} that have nothing to do with biology, and if you
ignore them a model may learn to \textbf{distinguish the batches rather than the
actual biological condition}. They are \emph{not} cured by simply collecting more
samples, and they are certainly \emph{not} beneficial ``variability'' --- they are a
confounder to be corrected (e.g.\ by batch-correction methods or careful
experimental design).

\section{Domain Shifts}

Medical models are deployed on data that drifts away from the training distribution.
Using $y$ for the disease status and $\xvec$ for patient features, the standard
taxonomy is:

\begin{itemize}
  \item \textbf{Covariate shift} --- the input distribution $p(\xvec)$ changes (so
        $p(\xvec)$ is certainly \emph{not} immune to shifts).
  \item \textbf{Label / prior shift} --- the class prevalence $p(y)$ changes; e.g.\
        during a COVID-19 wave the prevalence of the disease, and thus the
        probability of observing that phenotype, goes up.
  \item \textbf{Concept shift} --- the relationship $p(y\mid \xvec)$ changes; e.g.\
        if the diagnostic procedure that \emph{defines} $y$ is changed (an antigen
        test replaced by a PCR test), the label's meaning shifts.
\end{itemize}

\begin{notebox}[title={Validation accuracy does not buy robustness}]
High predictive quality on a randomly drawn validation set says \emph{nothing} about
whether a model will cope with future domain shifts --- it will \emph{not}
automatically adapt and maintain quality over time. Monitoring and re-validation are
required.
\end{notebox}

\section{Medical Imaging}

\subsection{Modalities}

Different physics, different images: \textbf{X-rays} pass radiation through the body
and differentiate tissues by \emph{density} (bone vs.\ soft tissue);
\textbf{Computed Tomography (CT)} is X-ray-based (\emph{not} ultrasound), stacking
many X-ray projections into a 3D image; \textbf{MRI} uses strong magnetic fields and
radio waves; \emph{ultrasound} uses sound waves; and \emph{optical coherence
tomography} uses \emph{light} (not X-rays).

\subsection{Tasks and the U-Net}

Compared with everyday computer vision, the dominant medical-imaging tasks are
\textbf{classification} and especially \textbf{segmentation} (delineating organs,
lesions, tumours pixel by pixel). The signature architecture is the
\textbf{U-Net}: a \emph{convolutional} network with a \textbf{symmetric encoder and
decoder} and, crucially, \textbf{skip connections} that concatenate encoder feature
maps into the corresponding up-sampling level of the decoder, preserving fine
spatial detail. It was built for biomedical \emph{segmentation} --- \emph{not} object
detection, and its encoder is convolutional, \emph{not} self-attention-based.

\subsection{Pitfalls and Interpretability}

\begin{notebox}[title={The ``Clever Hans'' effect}]
A model can appear to perform well while actually relying on \emph{irrelevant
features or artifacts} in the data (a ruler in tumour photos, a hospital tag) ---
shortcut learning. Good test scores can hide it, which is why interpretability
matters.
\end{notebox}

Explainable-AI tools help check \emph{why} a model decided what it did:
\textbf{Grad-CAM} highlights which image regions drove the prediction;
\emph{layer-wise relevance propagation} pushes a relevance signal back through the
network layer by layer; \emph{occlusion analysis} masks patches of the \textbf{input}
image and watches the prediction change (it perturbs inputs, not the training set);
and \emph{counterfactual explanations} generate an altered image showing what would
have to change for a different diagnosis.

\section{AlphaFold 2}

A landmark where life-science AI met structural biology: \textbf{AlphaFold 2}
predicts the \textbf{three-dimensional structure of a protein from its amino-acid
sequence}. It improves accuracy by incorporating \textbf{multiple sequence
alignments} and \emph{homologous} (template) structures, and its architecture is
\emph{attention}-based (the Evoformer) --- \textbf{not} a convolutional network, and
it uses \emph{no} quantum-mechanical calculations.


% ==========================================================
%  CHAPTER 8 — BIOLOGICAL SEQUENCES
% ==========================================================
\chapter{Biological Sequences}
\label{chap:sequences}

The last building block of life-science data is the \emph{sequence}. DNA, RNA and
proteins are all linear chains, which makes them a natural fit for the same
sequence-modelling toolbox used in NLP.

\section{What the Sequences Are Made Of}

\begin{itemize}
  \item \textbf{DNA} and \textbf{RNA} are chains of \emph{nucleic acids}
        (nucleotides).
  \item \textbf{Proteins} are chains of \emph{amino acids} linked by \emph{peptide
        bonds}.
\end{itemize}

\noindent These compositions are not interchangeable --- RNA is nucleic acid, not
lipid --- and biological sequences very much \emph{do} follow patterns and rules
(codons, motifs, secondary-structure propensities); they are not random strings.

\section{Deep Learning for Sequences}

To feed a sequence to a network it is typically \textbf{one-hot encoded} (each
nucleotide or amino acid a basis vector). From there, learned \textbf{feature
extraction} can identify \emph{motifs} --- recurring subsequences with biological
meaning (binding sites, domains). The full deep-learning toolbox applies to
\textbf{classification, segmentation, and generative modelling} of sequences, using
\textbf{1D-convolutions, RNNs, LSTMs, and Transformers} --- exactly the architectures
that also showed up for SMILES strings in Chapter~\ref{chap:generative}, which is no
coincidence: a SMILES string is just another biological-chemical sequence.

\begin{intbox}
Notice the unifying thread of the whole course: whether the object is a small
molecule (as a graph, fingerprint, or SMILES string), a protein (as a sequence or a
3D structure), an image, or a patient record, the recipe is the same --- choose a
representation that respects the object's structure and symmetries, then learn
$f_{\params}$ from data that is invariably limited, biased and noisy. Everything from
QSAR to AlphaFold is a variation on that single theme.
\end{intbox}



% ==========================================================
%  LIST OF ABBREVIATIONS
% ==========================================================
\chapter*{List of Abbreviations}
\addcontentsline{toc}{chapter}{List of Abbreviations}
\vspace{-1.0 cm}

\noindent {\small Branded model and database names used in the text (e.g.\ ChEMBL, ZINC, PubChem, REINVENT, DrugCLIP) are identified where they appear and are omitted here, as they are proper names rather than expandable abbreviations.}

\begin{longtable}{@{}p{0.16\textwidth}p{0.78\textwidth}@{}}
\toprule
\textbf{Abbreviation} & \textbf{Expansion} \\
\midrule
\endfirsthead
\toprule
\textbf{Abbreviation} & \textbf{Expansion} \\
\midrule
\endhead
\midrule
\multicolumn{2}{r}{\footnotesize\itshape continued on next page}\\
\endfoot
\bottomrule
\endlastfoot
AI & Artificial Intelligence \\
ADME & Absorption, Distribution, Metabolism, Excretion (pharmacokinetics) \\
ADMET & ADME plus Toxicity \\
AUC & Area Under the Curve \\
AUC-PR & Area Under the Precision--Recall curve \\
AUC-ROC & Area Under the ROC curve \\
BEDROC & Boltzmann-Enhanced Discrimination of ROC \\
CADD & Computer-Aided Drug Design \\
CAS & Chemical Abstracts Service \\
CASP & Computer-Assisted Synthesis Planning \\
CC$_{50}$ & Half-maximal Cytotoxic Concentration (50\% cell death) \\
CLIP & Contrastive Language--Image Pre-training \\
CMC & Chemistry, Manufacturing and Controls \\
CNN & Convolutional Neural Network \\
CT & Computed Tomography (X-ray-based 3D imaging) \\
CV & Cross-Validation \\
CYP450 & Cytochrome P450 \\
D-MPNN & Directed Message-Passing Neural Network \\
DMTA & Design--Make--Test--Analyse (cycle) \\
DNA & Deoxyribonucleic Acid (chain of nucleotides) \\
DNN & Deep Neural Network \\
ECE & Expected Calibration Error \\
ECFP & Extended-Connectivity Fingerprint \\
EF & Enrichment Factor \\
EGNN & E($n$)-Equivariant Graph Neural Network \\
EHR & Electronic Health Record \\
ESM & Evolutionary Scale Modeling (protein language model) \\
FCD & Fr\'echet ChemNet Distance \\
FNN & Feed-forward Neural Network \\
GAN & Generative Adversarial Network \\
GCN & Graph Convolutional Network \\
GFlowNet & Generative Flow Network \\
GNN & Graph Neural Network \\
GPT & Generative Pre-trained Transformer \\
Grad-CAM & Gradient-weighted Class Activation Mapping \\
hERG & human Ether-\`a-go-go-Related Gene (cardiac ion channel) \\
HTS & High-Throughput Screening \\
IC$_{50}$ & Half-maximal Inhibitory Concentration \\
InChI & International Chemical Identifier \\
IoU & Intersection over Union (Jaccard / Tanimoto) \\
IP & Intellectual Property \\
IUPAC & International Union of Pure and Applied Chemistry \\
$K_d$ & Dissociation Constant \\
$K_i$ & Inhibition Constant \\
KGCNN & Keras Graph Convolutional Neural Network (library) \\
KL & Kullback--Leibler (divergence) \\
LBDD & Ligand-Based Drug Design \\
LD$_{50}$ & Median Lethal Dose \\
LLM & Large Language Model \\
$\log P$ & octanol--water Partition coefficient (log scale) \\
LRP & Layer-wise Relevance Propagation \\
LSTM & Long Short-Term Memory \\
MACCS & Molecular ACCess System (keys) \\
MAE & Mean Absolute Error \\
MAML & Model-Agnostic Meta-Learning \\
MAPK & Mitogen-Activated Protein Kinase \\
MCC & Matthews Correlation Coefficient \\
MHN & Modern Hopfield Network \\
MIL & Multiple Instance Learning \\
ML & Machine Learning \\
MLP & Multilayer Perceptron \\
MPNN & Message-Passing Neural Network \\
MRI & Magnetic Resonance Imaging \\
MSA & Multiple Sequence Alignment \\
NN & Neural Network \\
OCT & Optical Coherence Tomography (light-based imaging) \\
ORD & Open Reaction Database \\
PAINS & Pan-Assay Interference Compounds \\
PCM & Proteochemometrics \\
PD & Pharmacodynamics (what the drug does to the body) \\
PDB & Protein Data Bank \\
PK & Pharmacokinetics (what the body does to the drug) \\
PPV & Positive Predictive Value (precision) \\
PROTAC & Proteolysis-Targeting Chimera \\
QED & Quantitative Estimate of Drug-likeness \\
QM & Quantum Mechanics / Mechanical \\
QSAR & Quantitative Structure--Activity Relationship \\
QSPR & Quantitative Structure--Property Relationship \\
$R^2$ & Coefficient of Determination \\
RF & Random Forest \\
RL & Reinforcement Learning \\
RMSE & Root Mean Squared Error \\
RNN & Recurrent Neural Network \\
Ro5 & (Lipinski's) Rule of Five \\
ROC & Receiver Operating Characteristic \\
SA & Synthetic Accessibility \\
SAR & Structure--Activity Relationship \\
SBDD & Structure-Based Drug Design \\
SELFIES & SELF-referencing Embedded Strings (always-valid line notation) \\
SMARTS & SMILES Arbitrary Target Specification \\
SMILES & Simplified Molecular Input Line Entry System \\
SVM & Support Vector Machine \\
TFN & Tensor Field Network \\
TNR & True Negative Rate (specificity) \\
TPR & True Positive Rate (recall / sensitivity) \\
TPSA & Topological Polar Surface Area \\
USPTO & United States Patent and Trademark Office \\
VAE & Variational Autoencoder \\
VS & Virtual Screening \\
XAI & Explainable Artificial Intelligence \\
XC$_{50}$ & Half-maximal Effect Concentration (the IC$_{50}$/EC$_{50}$/CC$_{50}$ family) \\
\end{longtable}

\end{document}